% Modernised reading — fonds Alexandre Grothendieck, shelfmark 70
% (pages 1-125, the whole folder).
%
% This is an INTERPRETATION, not a transcription. It re-reads the pages in
% today's notation and vocabulary, states the hypotheses the manuscript leaves
% implicit, and drops the critical apparatus so that the mathematics reads
% straight through. Everything the brackets and underlines carried moves into
% a footnote. It is held to being mathematically correct as it stands; where
% the manuscript is loose or mistaken, the true statement is what appears here
% and the footnote says what the page has. In any dispute of substance the
% transcription governs: whoever cites Grothendieck must cite batch-01.fr.tex
% (pp. 1-20), batch-02.fr.tex (pp. 21-40), batch-03.fr.tex (pp. 41-60),
% batch-04.fr.tex (pp. 61-80), batch-05.fr.tex (pp. 81-100), batch-06.fr.tex
% (pp. 101-120) or batch-07.fr.tex (pp. 121-125).
%
% Pass: Opus 5.5 (claude-opus-5-5), 2026-10-10 — first pass, unchecked against the pages by a human.
% The folder was read whole, its seven batches together; this is one document
% for the shelfmark. It was made on Opus 5.5 and has not been measured against
% the reference reading of folder 115 (made on Opus 5). The literature named
% in the footnotes is cited from memory and was not checked against the
% sources. Checked by machine (sympy, exact arithmetic; finite fields by hand-
% written code): the invariance of the quadratic function f of p. 68 under
% sigma_0, sigma_1, sigma_2 for general alpha, beta; the discriminants delta'
% and delta of p. 68 and the table of p. 69; the centre and f(centre) of
% pp. 68-69; the relations and the characteristic polynomials of
% sigma_{i-1} sigma_i for the normal form of pp. 18-25 (n = 2, 3, 4); for the
% icosahedron over Q(sqrt 5): the group has order 120, the orbit of s_0 is the
% twelve points listed, (sigma_0 sigma_1 sigma_2)^5 is x -> 2o - x, the
% quadrics through the twelve vertices form a line (rank 9), D_2 of p. 111;
% over Z[alpha]/(alpha^2+alpha-1): the twelve points are distinct, no three
% are collinear and no vertex lies in the plane of a face not containing it,
% in EVERY characteristic (all the relevant minors are units of Z[alpha]);
% Delta = alpha Delta' and Phi' = (2 alpha - 1) Phi of p. 76.
%
% Scope. Covers (no \pagerange of their own; their words head the parts):
% p. 1 « n-polyèdres et n-hyperpolyèdres », p. 2 « n-polyèdres réguliers en
% coordonnées affines universelles (vieilles rédactions) (23 p) » [= pp. 3-25],
% p. 28 « n-quasi-polyèdres (vieilles rédactions) », p. 38 « “étude locale” des
% n-polyèdres et des n-cartes (vieilles rédactions) », p. 49 « 2 – Polyèdres
% réguliers. (Vieille rédaction) », p. 50 « (18) 2-Polyèdres repérés en
% coordonnées universelles (vieille rédaction) », p. 73 « (32) Calculs en
% coordonnées universelles (icosaèdre) — vieilles rédactions », p. 113
% « Position du pb des 2-polyèdres réguliers / vieille rédaction (10 p) ».
% No mathematics: 26, 27, 37, 48, 60, 63, 72, 112, 125 (blank or separators),
% 57 (a draft letter in his hand, cancelled — not described), 77, 79, 124
% (typed versos; 124 a USTL letter stamped 17 June 1976 per the transcription's
% header, not restated). Page 81 is not transcribed; the transcription does not
% say why. Pages 100 (homotopies), 103 and 105 (mu_n, Sym^2 F) are sheets of
% another subject or paper, described briefly.
%
% Spine.
%   pp. 3-25    regular n-polyhedra: universal regular polytope of type
%               {p_1..p_n} (string Coxeter group); regular affine realisation
%               over a ring; the frame of vertices s_0..s_{n+1}; parameters;
%               requiring reflections leaves n parameters alpha_0..alpha_{n-1};
%               normal form sigma_i = id + l_i (x) (s_{i+1} - s_i).
%   pp. 29-36   n-quasi-polyhedra: same frame without sigma_i^2 = 1;
%               2n+1 parameters; orders (n = 1).
%   pp. 39-47   local study: boundary and co-face of a facet, sections
%               [f_i, f_j], the invariants p_i; geometric representation by
%               maximal flags; induced realisations.
%   pp. 51-59   2-polyhedra realised by points, lines, planes; the four first
%               vertices are independent; regular realisation = (phi_G, a
%               flag); p. 56 which standard maps are incidence geometries /
%               polyhedra; p. 58 an isolated sheet on groupoids.
%   pp. 61-62   PSL(2,Z) and its quotients mod n (n <= 5: Platonic groups).
%   pp. 64-71   n = 2 in coordinates: alpha (vertex), beta (face); the
%               invariant quadratic function; table for the Platonic solids
%               and the pentagrammic ones; beta = 2: a hyperboloid.
%   pp. 74-111  the icosahedron in « universal coordinates » over
%               Z[alpha]/(alpha^2+alpha-1): reflections, rho_f^3, rho_s^5,
%               twelve vertices, centre and antipodism (char 2: a translation),
%               the circumscribed quadric (char 5: a cone), remarkable vectors,
%               regular and faithful realisations.
%   pp. 114-119 the problem of regular 2-polyhedra in a scheme X/S: universal
%               flags, axioms [11]-[20], [32]-[33'], N, Gamma/N, questions.
%   p. 120      « cocubes »: the cube and the cross-polytope from a set with
%               involution.
%   pp. 121-123 « Programme provisoire d'une théorie arithmétique des
%               polyèdres réguliers »: k-realisations of a real type tau in a
%               projective space; a sheet of counting notation.
%
% Notation fixed for the document.
%  - G_n (written \mathfrak{G}_n, his gothic G) = <sigma_0..sigma_n | sigma_i^2,
%    (sigma_i sigma_j)^2 for |i-j| >= 2>, the cartographic group with n+1
%    generators (folder 76's G_n; for n = 2 folder 88's hat-Gamma, which p. 114
%    writes Gamma). sigma_i changes the i-dimensional member of a flag.
%  - N the stabiliser of a flag (p. 3 writes H), Gamma_{p_*} the Coxeter
%    group [p_1..p_n], p_i = order of sigma_{i-1} sigma_i (p. 3's indexing;
%    p. 42 indexes from 0). For J in {0..n}: G_J = <sigma_i, i not in J>
%    (p. 4 writes G_{n,H}), D_J flags of type J.
%  - s_0 .. s_{n+1} the frame of vertices, e_i = s_i - s_0, x_1..x_{n+1}
%    affine coordinates, Sigma_i = x_{i+2} + ... + x_{n+1} (Sigma_0 =
%    x_2 + ... + x_{n+1}); alpha_0..alpha_{n-1} as on pp. 15-25.
%  - n = 2: beta = alpha_0 (face), alpha = alpha_1 (vertex), as on pp. 64-111;
%    rho_f = sigma_0 sigma_1 (about a face, order q), rho_s = sigma_1 sigma_2
%    (about a vertex, order p), as on p. 67 and in folder 88 (p. 61 swaps the
%    names; p. 84 writes u, v). alpha' = -1 - alpha.
%  - Psi_m the minimal polynomial of 2 cos(2 pi / m), his « polynôme
%    semi-cyclotomique » F_m.
%  - Icosahedron: s'_3 = sigma_1 s_3, s''_3 = rho_f s_3 (pp. 74-111; p. 64
%    names sigma_1 s_3 s''_3 and sigma_0 s_2, sigma_0 s_3 s'_2, s'_3, here
%    t_2, t_3). iota the antipodism (his underlined a, as in folder 86), o the
%    centre, q the quadratic function, q_0 its homogeneous part.
%  - « réflexion »: an affine involution whose fixed locus is, in every fibre,
%    an affine hyperplane (projective: a projective hyperplane), as the pages
%    use the word; in characteristic 2 such a map is a transvection.
%  - The base scheme of pp. 110 and 115 is written Y (the pages write S, the
%    letter of the vertex set).
%
% Substantive departures from the pages, each footnoted where it occurs:
%  p. 4    R_{p_*} = G_n / Gamma_{p_*} read G_n / N (= Gamma_{p_*} as a set).
%  p. 9    sigma_1 s_4 in f_4 \ f_2 read f_4 \ f_3.
%  p. 11   the boxed sigma_i s_j = s_0 + lambda_i (e_i + e_{i+1}) + e_j has the
%          wrong sign: the involution condition forces lambda_{ii} =
%          -lambda_{i,i+1}; the form is s_0 + lambda_i (e_{i+1} - e_i) + e_j
%          (as pp. 18 and 25 have).
%  pp. 13-18 the reflection requirement is his addition to the definition;
%          sigma_i^2 = 1 alone allows lambda_{i,i+2} = -1. Over a ring the
%          page argues through nilpotents and residue characteristic 2; the
%          comparison of coefficients given here is ours.
%  p. 18   x_{n+2} read x_{n+1}; « lambda_{i,i+1} + lambda_{i,i+1} » read
%          lambda_{ii} + lambda_{i,i+1}.
%  p. 19   the commutation of sigma_i, sigma_j (j >= i+2) needs that sigma_j
%          preserves x_j + x_{j+1}; supplied.
%  pp. 19-22 the orders of sigma_{i-1} sigma_i are not finished on the page;
%          the characteristic polynomial (t-1)^{n-1}(t^2 - alpha_{i-1} t + 1)
%          and the criterion « alpha_{i-1} = zeta + zeta^{-1}, zeta primitive
%          p_i-th root » are ours (machine-checked for n <= 4, proved in
%          general).
%  p. 33   X_i lacks the term (1 - lambda_i) mu_i Sigma_i and X_{i+1} has
%          coefficient lambda_i^2 - lambda_i + 1 for x_{i+1} (the page: 1);
%          the conclusions (21) and (22) stand.
%  p. 34   cases a), b) of sigma_{0v}^n = 1 corrected: a) 1 + eta + ... +
%          eta^{n-1} = 0, i.e. (eta != 1, eta^n = 1) or (eta = 1, p | n);
%          b) eta = 1, alpha_0 = 0, p does not divide n.
%  p. 35   « alpha_0 = 1: sigma_0 = id » is impossible (sigma_0 s_0 = s_1);
%          for eta_0 = 1 sigma_0 is (x,y) -> (x + 1 + alpha_0 y, y), of order p
%          (p > 0) or infinite, whatever alpha_0.
%  p. 36   the alpha_i of p. 36 is not that of p. 18 (1 + alpha_i of p. 36 is
%          -(1 + alpha_i) of p. 18); eigenvalue 1 has multiplicity n, not n-1.
%  p. 39   the signs ± on Gamma_n(i), E(i) read as the definitions require
%          (quotient by the upper group, the lower one acts).
%  p. 42   « true at least for incidence geometries »: the equality p_k =
%          order of sigma_k sigma_{k+1} holds under the intersection property.
%  p. 45   Drap aff(Pi) read Drap aff(E).
%  p. 52   « phi_S(S) generates E » needs dim E = 3; supplied.
%  p. 54   sigma_1(a_0) = a_0 read sigma_1(f_0) = f_0.
%  p. 56   the bracketed statements (which correct a) and b)) are followed.
%  p. 58   « centralisateur de C » read « de G ».
%  p. 61   rho'_s = sigma_2 rho_f sigma_2^{-1} read sigma_2 rho_s sigma_2^{-1}
%          (in p. 61's names).
%  p. 64   s''_3 = s_3 + xi e_1 + eta e_2 + zeta e_3 has third coordinate
%          1 + zeta, not 1 - zeta; then sigma_1^2 = 1 gives zeta(zeta + 2) = 0,
%          not zeta^2 = 0; with det = -1 (or a fixed hyperplane) zeta = 0, as
%          the page concludes.
%  p. 66   sigma_0 is a reflection iff mu = 1; for mu = -1 it is the central
%          symmetry about the midpoint of s_0 s_1.
%  p. 68   the doubtful factor (alpha - beta) of delta' is (alpha + beta)
%          (checked).
%  p. 74   « H_2 parallel to H_0 in char. 2 » holds on the equations as
%          transcribed (the transcription's note doubts it).
%  p. 75   iota s''_3 = (alpha', 1 - 2 alpha', -alpha') is (alpha', 1,
%          -alpha') (checked); iota s'_3 = (0, 1, -alpha') is right, and so
%          the expectation in the note at p. 92 is not.
%  p. 84   the second matrix of sigma_1 has alpha', -alpha' in its third
%          column; the cycle of rho_f on s_0 ends at s_0.
%  p. 85   sigma_0 sigma_2 = sigma_2 sigma_1 read sigma_2 sigma_0.
%  p. 91   in char. 2 iota is the translation by (alpha', alpha, alpha)
%          = (alpha', 1 - alpha', 1 - alpha') (the page reads (alpha', 1, 1)).
%  p. 93   « four distinct vertices are coplanar only if on a face » is false
%          in every characteristic (the five neighbours of a vertex are
%          coplanar — p. 85 says so for four of them —, and so are two opposite
%          edges); the restricted form the page then proves suffices, and
%          holds in every characteristic (checked).
%  p. 102  the third row of the 4x4 matrix is (1, -alpha, 2, -1); its
%          determinant is -alpha - 3 (checked), not the alpha' of p. 104.
%  p. 108  lambda^2 + lambda - 1 = 0 does not express rho_f^3 = id (which is
%          automatic, p. 85) but rho_s^5 = id.
%  p. 117  (23)-(24) define an ANTI-isomorphism Gamma_P -> Aut(X, P).
%  p. 123  N = prod s_i over 0 <= i <= n (the bounds read 1, n-1 do not give
%          card G); the meaning of s_i, s'_i is our reading.
%
% Links the pages do not write (ours, said so in the body): the realisations
% of pp. 3-25 are reflection representations of the Coxeter group [p_1..p_n]
% (for a finite group: Tits's geometric representation or a Galois conjugate);
% the « semi-cyclotomic » condition of p. 67 is the c of p. 103 (Phi_n(c) = 0,
% c = zeta + zeta^{-1}); beta = 2 on p. 70 is the hyperboloid model of the
% hyperbolic plane carrying the tessellation {infinity, p}; pp. 74-111 carry
% out the « icosaèdre régulier général sur un anneau fini sur Z » announced by
% the course of folder 75 (1977-78); the faithfulness asked on p. 21 and
% p. 109 holds for the icosahedron in every characteristic; p. 105's
% Sym^2(F) (x) (det F)^{-1} with its discriminant is folder 82's PGL_2 = SO_3;
% p. 120's cube is folder 76's cube of a set with involution, and its
% « cocube » the cross-polytope of folder 87; pp. 114-119 are the
% scheme-theoretic form of folder 86's part II.
%
% Announced and never established: the computation of (sigma_0 sigma_1)^p
% (p. 19) and the « ? » of p. 20; whether D_i -> Grass_i is injective (p. 21);
% the realisation of the co-face (p. 47, breaks off); the « à vérifier » of
% pp. 53 and 56; the meaning of the exact sequences of p. 58; v^5 = id <=>
% alpha^2 + alpha - 1 = 0 (p. 85, « un calcul commode »; true, checked); that
% G has order 120 (p. 87, « si nous admettons »; true, classical); the
% Proposition « régulière => fidèle » (p. 109, proof breaks off); the
% classification programme in the margin of p. 111; questions (25)-(29)
% (pp. 117-118); conditions 1)-4) of p. 122 are a definition with no result.
\documentclass[11pt,a4paper]{article}
\input{../preamble/grothendieck}

\folder{70}
\pages{1}{125}
\dating{[à partir de 1976]}
\watermark{Édition de démonstration\\interprétation personnelle de l'œuvre}
\foldertitle{Réalisations géométriques de structures combinatoires (n-polyèdres, n-hyperpolyèdres [2-polyèdres réguliers]…) (Vieilles rédactions) : notes manuscrites (s.d.) — lecture modernisée du dossier entier}

\begin{document}

\section*{Résumé}

\begin{resume}
Le cube, l'octaèdre, l'icosaèdre sont des polyèdres « réguliers » : toutes
leurs faces sont pareilles, tous leurs sommets aussi, et ils ont autant de
symétries qu'on peut en avoir. Les Grecs les définissaient par des longueurs et
des angles égaux. Ces cent vingt-cinq pages partent de l'autre bout. On se
donne d'abord le polyèdre sans aucune géométrie --- une liste de sommets,
d'arêtes et de faces, qui touche qui, et le groupe de ses symétries --- et l'on
demande de combien de façons on peut le \emph{réaliser} : le dessiner dans un
espace par des points, des droites et des plans, de façon que toutes ses
symétries combinatoires deviennent des transformations de l'espace. Et l'espace
n'est pas forcément celui des nombres réels : ce peut être un espace sur un
corps fini, ou sur un anneau quelconque.

L'image qui guide tout est celle du kaléidoscope. Trois miroirs bien placés
et un petit motif suffisent à fabriquer une figure symétrique entière ; ici
le motif est un \emph{drapeau} --- un sommet, une arête qui en part, une face
que borde l'arête --- et les miroirs sont trois symétries élémentaires, qui
changent l'un des trois éléments du drapeau en gardant les deux autres. Pour un
polyèdre de dimension supérieure, il y a un miroir de plus par dimension.
Grothendieck place les sommets du premier drapeau aux points les plus simples
possibles --- l'origine et les extrémités des vecteurs de base --- et calcule
les miroirs dans ces « coordonnées universelles ». La surprise est que tout se
réduit à une poignée de nombres, un par rotation élémentaire du polyèdre, et
que chacun est imposé par la forme des faces : pour une face triangulaire, le
nombre vaut $-1$ ; pour un pentagone, il doit vérifier $\alpha^2 + \alpha = 1$,
l'équation du nombre d'or. Ces équations sont à coefficients entiers, et ont
donc un sens en toute caractéristique.

Le cœur du dossier est l'icosaèdre ainsi traité. On le voit exister sur tout
anneau où l'équation du nombre d'or a une solution, et l'on suit ce qu'il
devient en chemin. En caractéristique $2$, le centre de l'icosaèdre part à
l'infini et la symétrie centrale devient une translation. En caractéristique
$5$, la « sphère » circonscrite dégénère en un cône : l'icosaèdre est inscrit
dans une sphère de rayon nul. Pourtant, en aucune caractéristique trois
sommets ne deviennent alignés, et aucun sommet ne tombe dans le plan d'une face
qui ne le contient pas : l'icosaèdre reste partout un icosaèdre. Chemin
faisant, une forme quadratique apparaît dont le discriminant vaut exactement
$1$, et la page s'exclame « merveille ! » ; plus loin, une vérification qui
aboutit est saluée d'un « et on est heureux ! ».

Le dossier porte, sur ses chemises, la mention « vieilles rédactions », et
c'en est bien : la même question est reprise six ou sept fois, chaque fois
d'un autre côté --- en dimension quelconque, puis sans supposer que les
miroirs soient des symétries, puis par l'étude locale des faces, puis en
dimension trois, puis sur une base quelconque, enfin, dans un « programme
provisoire d'une théorie arithmétique des polyèdres réguliers », en partant du
modèle réel et en lui cherchant des réalisations sur tout anneau. On y voit
une façon de travailler : refaire le même calcul jusqu'à ce qu'il ne dépende
plus d'aucun choix, et poser la question dans le cadre le plus large où elle a
encore un sens. Un aparté rattache les groupes des polyèdres de Platon au
groupe modulaire $\mathrm{PSL}(2, \mathbf{Z})$, réduit modulo $2, 3, 4, 5$.

Les noms d'aujourd'hui sont ceux de polytope abstrait régulier, de groupe de
Coxeter et de sa représentation géométrique, de groupe cartographique, de
réalisation d'un polytope, de sous-groupe de congruence, et, pour la dernière
page, de polytope croisé.

\keywords{abstract regular polytope, string Coxeter group, intersection property, Schläfli symbol, cartographic group, flag, realization of polytopes, affine reflection, transvection, Tits geometric representation, minimal polynomial of 2cos(2π/n), invariant quadratic form, discriminant, modular group, principal congruence subgroup, icosahedron, golden ratio, characteristic 2, characteristic 5, Kepler-Poinsot polyhedra, hyperboloid model, regular tessellation of the hyperbolic plane, 1-cocycle, scheme-theoretic polyhedron, cross-polytope}
\end{resume}

\pagerange{1}{123}
\section*{Le fil du dossier, et les conventions}

\subsection*{Les chemises et les stations}

Huit feuilles de chemise, de sa main, découpent le dossier ; quatre portent un
nombre de pages, et l'un au moins se vérifie exactement (« (23 p) » pour les
pages 3 à 25). Elles donnent ses titres aux parties de cette lecture. Le mot
« $n$-hyperpolyèdres » de la première chemise n'est repris nulle part ailleurs.

\begin{itemize}
  \item \textbf{Pages 1 à 25 : « $n$-polyèdres réguliers en coordonnées affines
    universelles ».} Le polyèdre régulier universel de type $\{p_1, \dots,
    p_n\}$ ; ses réalisations affines régulières sur un anneau ; le repère des
    sommets ; l'exigence que les générateurs soient des réflexions laisse $n$
    paramètres ; une forme normale.
  \item \textbf{Pages 28 à 36 : « $n$-quasi-polyèdres ».} Le même repère sans
    supposer que les générateurs soient des involutions : $2n+1$ paramètres.
  \item \textbf{Pages 38 à 47 : « étude locale » des $n$-polyèdres et des
    $n$-cartes.} Bord et étoile d'une facette, sections, invariants ;
    représentation géométrique par des drapeaux, réalisations induites.
  \item \textbf{Pages 49 à 59 : « 2-polyèdres repérés en coordonnées
    universelles ».} Le cas de la dimension trois, avec ses points, droites et
    plans ; un feuillet isolé sur les groupoïdes (page 58).
  \item \textbf{Pages 61 et 62 : le groupe modulaire} et ses quotients
    modulo $n$.
  \item \textbf{Pages 64 à 71 : les trois réflexions en coordonnées, et la
    quadrique invariante}, pour tous les types à la fois ; le tableau des
    solides de Platon et des polyèdres étoilés ; un cas hyperbolique.
  \item \textbf{Pages 73 à 111 : « calculs en coordonnées universelles
    (icosaèdre) ».} Trois passes sur l'icosaèdre au-dessus de l'anneau
    $\mathbf{Z}[\alpha]/(\alpha^2 + \alpha - 1)$ : les réflexions, les douze
    sommets, le centre et l'antipodisme, la quadrique circonscrite et ses
    discriminants, les réalisations régulières et fidèles. Trois feuilles
    intercalées traitent d'autre chose (pages 100, 103, 105).
  \item \textbf{Pages 113 à 119 : « position du problème des 2-polyèdres
    réguliers ».} Drapeaux universels, polyèdre régulier dans un schéma,
    axiomes, questions.
  \item \textbf{Page 120 : les « cocubes ».}
  \item \textbf{Pages 121 à 123 : « programme provisoire d'une théorie
    arithmétique des polyèdres réguliers ».}
\end{itemize}

Aucune de ces rédactions ne renvoie à une autre. Elles sont rangées ici dans
l'ordre des feuillets, qui n'est pas forcément celui de leur écriture : la
dernière chemise porte « (10 p) », ce qui la ferait courir jusqu'à la
page 123, mais les pages 120 à 123 ont leurs propres sujets.

\textbf{Les dossiers voisins.} Le groupe cartographique à $n+1$ générateurs
est celui de la lecture du dossier 76, qui le note $G_n$ ; pour $n = 2$, c'est
le $\widehat{\Gamma}$ des dossiers 86 et 88, dont on reprend les rotations
$\rho_f$, $\rho_s$ et les cartes standards $C_{p,q}$. Les drapeaux, les trois
involutions et l'antipodisme de l'icosaèdre sont ceux du dossier 86, qui
traite le même icosaèdre sans coordonnées. Les annonces de cours de
1977--1978 du dossier 75 décrivent exactement ce que font les pages 74 à
111 : réaliser l'icosaèdre « sur les réels et sur d'autres corps », la
caractéristique $2$ et la caractéristique $5$ « jouant un rôle particulier »,
et construire « l'icosaèdre régulier général sur un anneau de base fini
sur $\mathbf{Z}$ ». La forme quadratique de la page 105 est celle de
l'isomorphisme $\mathrm{PGL}_2 \simeq \mathrm{SO}_3$ du dossier 82 ; le cube
de la page 120 est celui du dossier 76, et le « cocube » le polytope croisé du
dossier 87.

\subsection*{Conventions, valables pour tout le dossier}

\textbf{Le groupe cartographique.} Pour $n \geq 0$,
\[
\mathfrak{G}_n = \bigl\langle \sigma_0, \dots, \sigma_n \bigm| \sigma_i^2 = 1,\
(\sigma_i \sigma_j)^2 = 1 \text{ si } |i - j| \geq 2 \bigr\rangle ,
\]
noté par lui d'un $G$ gothique (page 114 : $\Gamma$, pour $n = 2$). Un
$\mathfrak{G}_n$-ensemble est une \emph{$n$-carte} ; ses éléments sont les
\emph{repères}, ou drapeaux maximaux $(f_0, f_1, \dots, f_n)$, et $\sigma_i$
change la composante $f_i$ en gardant les autres. Une $n$-carte a des facettes
de dimensions $0$ à $n$ ; ses réalisations vivent dans un espace de dimension
$n + 1$, dont elle est, en quelque sorte, le bord : les « $2$-polyèdres » sont
les polyèdres usuels de l'espace à trois dimensions.

\textbf{Types.} $p_i$ est l'ordre de $\sigma_{i-1}\sigma_i$ ($1 \leq i \leq n$),
comme à la page 3 ; $\{p_1, \dots, p_n\}$ est alors le symbole de Schläfli.
Pour $n = 2$ on suit le dossier 88 : $\rho_f = \sigma_0\sigma_1$ tourne autour
d'une face, d'ordre $q = p_1$ (le nombre de côtés), $\rho_s = \sigma_1\sigma_2$
tourne autour d'un sommet, d'ordre $p = p_2$ (le nombre d'arêtes en un
sommet)\footnote{La page 67 nomme ainsi $\rho_s$ et $\rho_f$ ; la page 61 les
échange, et la page 84 les appelle $u$ et $v$. L'icosaèdre, de symbole
$\{3, 5\}$, a donc $q = 3$ et $p = 5$.}. $\Psi_m$ désigne le polynôme minimal
de $2\cos(2\pi/m)$ sur $\mathbf{Q}$, c'est-à-dire de $\zeta + \zeta^{-1}$ pour
$\zeta$ racine primitive $m$-ième de l'unité : c'est son « polynôme
semi-cyclotomique » $F_m$ (page 67). Ainsi $\Psi_3(\alpha) = \alpha + 1$,
$\Psi_4(\alpha) = \alpha$, $\Psi_5(\alpha) = \alpha^2 + \alpha - 1$.

\textbf{Le repère des sommets.} Un drapeau $f_0 \subset \cdots \subset f_n$ de
sous-espaces affines étant donné, $s_0 = f_0$ et $s_{i+1} = \sigma_i s_i$ ;
$e_i = s_i - s_0$ ; $(x_1, \dots, x_{n+1})$ sont les coordonnées affines dans le
repère $(s_0; e_1, \dots, e_{n+1})$, et
\[
\Sigma_0 = x_2 + \cdots + x_{n+1}, \qquad
\Sigma_i = x_{i+2} + \cdots + x_{n+1} \quad (1 \leq i \leq n-1) .
\]
Les paramètres $\alpha_0, \dots, \alpha_{n-1}$ sont ceux des pages 15 à 25.
Pour $n = 2$, les pages 64 à 111 écrivent $\beta = \alpha_0$ (la face) et
$\alpha = \alpha_1$ (le sommet) ; on les suit. Pour l'icosaèdre, $\alpha' = -1
- \alpha$ est la racine conjuguée de $\alpha^2 + \alpha - 1$.

\textbf{L'icosaèdre.} $s'_3 = \sigma_1 s_3$ et $s''_3 = \rho_f s_3$, comme
aux pages 74 à 111\footnote{La page 64 appelle $s''_3$ le point $\sigma_1 s_3$
et $s'_2$, $s'_3$ les points $\sigma_0 s_2$, $\sigma_0 s_3$ ; on les note ici
$s'_3$, $t_2$, $t_3$.} ; $\iota$ est l'antipodisme, qu'il note d'un $a$
souligné (on garde la lettre du dossier 86) ; $o$ est le centre ; $q$ la
fonction quadratique invariante et $q_0$ sa partie homogène.

\textbf{Réflexion.} Les pages appellent \emph{réflexion} une involution
affine dont l'ensemble des points fixes est, dans chaque fibre, un hyperplan
affine ; « réflexion projective » quand cet hyperplan peut être à l'infini. On
garde le mot. En caractéristique $2$, une telle involution est une
transvection.

\subsection*{Ce que seule une lecture d'ensemble peut dire}

\textbf{Le même calcul, refait.} La forme normale des générateurs est établie
trois fois : en dimension quelconque (pages 9 à 25), sans l'hypothèse
d'involution (pages 30 à 36), et en dimension trois (pages 64 à 66, puis 82 à
84) ; ces trois versions ne nomment pas leurs points ni leurs paramètres de la
même façon, et les notes le signalent où il le faut. Le paramètre
« invariant modulaire du polygone régulier » de la page 15, le
« polynôme semi-cyclotomique » de la page 67 et le $c = \zeta + \zeta^{-1}$
avec « $\Phi_n(c) = 0$ » de la page 103 sont une seule et même chose : la
trace de la rotation élémentaire.

\textbf{Une hypothèse ajoutée.} Que les générateurs soient des réflexions
n'est pas dans la définition d'une réalisation (pages 5 et 6, 43, 51) ; il
l'impose à chaque fois au moment du calcul (pages 13, 15, 31, 66, 83). Les
pages 29 à 36 sont la tentative de s'en passer.

\textbf{Une question, et sa réponse pour l'icosaèdre.} La question de la
page 21 --- une réalisation régulière envoie-t-elle des facettes distinctes sur
des sous-espaces distincts ? --- et la « fidélité » des pages 109 et 110 ont,
pour l'icosaèdre, une réponse positive en toute caractéristique : les pages
93 et 94 la démontrent en partie, et la vérification complète est faite dans
cette lecture.

\textbf{Deux points de vue sur le même objet.} Les pages 114 à 119 cherchent à
définir un polyèdre régulier dans un schéma par des axiomes, puis montrent
qu'il revient à un morphisme $\varphi : \mathfrak{G}_2 \to \mathrm{Aut}(X)$ et à un
drapeau ; c'est exactement la donnée que les pages 74 à 111 construisent pour
l'icosaèdre. Le programme des pages 121 et 122 renverse la perspective : il
part du modèle réel et lui cherche des réalisations sur un anneau.

\subsection*{Ce que le dossier annonce sans l'établir}

Le calcul de $(\sigma_0\sigma_1)^p$ (page 19) et les deux « $\overset{?}{=}$ »
de la page 20 restent ouverts (on les complète) ; la question d'injectivité
de la page 21 n'est pas tranchée en général ; la réalisation induite de
l'étoile d'une facette s'arrête à la page 47 ; deux « à vérifier » (pages 53
et 56) ; les suites exactes de la page 58 ; « un calcul commode » (page 85) et
« si nous admettons que $G$ est d'ordre $120$ » (page 87), tous deux exacts ;
la proposition « une représentation régulière est fidèle » (page 109), dont la
démonstration s'interrompt ; le programme de classification esquissé en marge
de la page 111 ; les questions (25) à (29) des pages 117 et 118.

\pagerange{1}{25}
\section*{I. Les $n$-polyèdres réguliers en coordonnées affines universelles (pages 1 à 25)}

La chemise de la page 1 porte « $n$-polyèdres et $n$-hyperpolyèdres » ; celle de
la page 2, « $n$-polyèdres réguliers en coordonnées affines universelles
(vieilles rédactions) » et « (23 p) » : ce sont les pages 3 à 25.

\pagerange{3}{5}
\subsection*{Le polyèdre régulier universel (pages 3 à 5)}

Soit $R$ une $n$-carte \emph{régulière} : un $\mathfrak{G}_n$-ensemble transitif
dont le stabilisateur $N$ d'un repère $r$ est distingué dans $\mathfrak{G}_n$ ---
il est alors le même pour tous les repères --- ; le groupe $\Gamma =
\mathfrak{G}_n / N$ opère simplement transitivement sur $R$\footnote{La page
appelle $H$ ce stabilisateur ; elle réemploie la lettre $H$ à la page 4 pour
une partie de $\{0, \dots, n\}$. On l'appelle ici $N$, et $J$ la partie.}. On
suppose que pour toute partie $I \subset \{0, \dots, n\}$ sans deux entiers
consécutifs, le produit $\sigma_I = \prod_{i \in I} \sigma_i$ est $\neq 1$ dans
$\Gamma$, et l'on pose $p_i = $ ordre de $\sigma_{i-1}\sigma_i$, avec $3 \leq p_i
\leq \infty$ (« les $p_i$ peuvent être infinis ! »). On se place dans le cas
où $\Gamma$ n'a pas d'autres relations que celles qu'impose ce qui précède :
\[
\Gamma = \Gamma_{p_*} = \bigl\langle \sigma_0, \dots, \sigma_n \bigm|
\sigma_i^2,\ (\sigma_i\sigma_j)^2 \ (|i-j| \geq 2),\ (\sigma_{i-1}\sigma_i)^{p_i}
\bigr\rangle .
\]
C'est le \emph{groupe de Coxeter} de type $[p_1, \dots, p_n]$, à diagramme
linéaire\footnote{Le nom n'est pas sur la page, ni celui de « symbole de
Schläfli » ; la page écrit $p_* = (p_1, \dots, p_n) \in ([3, +\infty] \cap
\mathbf{N})^n$.}. Dans un groupe de Coxeter, un produit de générateurs
distincts n'est jamais trivial, de sorte que les conditions $\sigma_I \neq 1$
sont satisfaites, comme l'affirme la page 4. Toutes les cartes $R$ de ce type
sont isomorphes, et canoniquement isomorphes une fois choisi un repère : $R
\simeq \Gamma_{p_*}$, le repère étant l'élément neutre\footnote{La page écrit
$\underline{R}_{p_*} = \mathfrak{G}_n / \Gamma_{p_*}$, ce qui n'a pas de sens
($\Gamma_{p_*}$ est un quotient de $\mathfrak{G}_n$, non un sous-groupe) ; on lit
$\mathfrak{G}_n / N$, qui est $\Gamma_{p_*}$ comme ensemble.}.

Pour $J \subset \{0, \dots, n\}$, soit $\mathfrak{G}_J$ le sous-groupe engendré
par les $\sigma_i$, $i \notin J$, et $D_J = R / \mathfrak{G}_J$ l'ensemble des
\emph{drapeaux de type $J$} ; $J \subset K$ entraîne $\mathfrak{G}_K \subset
\mathfrak{G}_J$, d'où $D_K \to D_J$, et $D_i = D_{\{i\}}$ est l'ensemble des
$i$-facettes. La page 5 affirme, « sauf erreur », que la carte ainsi définie
est une \emph{géométrie d'incidence} dont $D_i$ est l'ensemble des facettes de
dimension $i$, c'est-à-dire un polyèdre combinatoire régulier. C'est exact :
un groupe de Coxeter linéaire vérifie la « propriété d'intersection »
$\mathfrak{G}_J \cap \mathfrak{G}_K = \mathfrak{G}_{J \cup K}$ (dans $\Gamma$),
et $R$ est l'ensemble des drapeaux du \emph{polytope régulier universel}
$\{p_1, \dots, p_n\}$, dont tout polytope régulier de ce type est un
quotient\footnote{Propriété d'intersection et polytope universel : Tits pour
les groupes de Coxeter, McMullen et Schulte, \emph{Abstract Regular
Polytopes} (2002), pour la théorie des polytopes abstraits ; cité de mémoire.
La condition $p_i \geq 3$ écarte les cas dégénérés ; la page 56 y revient.}.

\pagerange{5}{6}
\subsection*{Réalisations affines régulières (pages 5 et 6)}

Soit $k$ un anneau. Une \emph{réalisation affine régulière} de $\Pi$ (le
polyèdre défini par $R$) est la donnée
\begin{enumerate}
  \item[a)] d'un espace affine $E$ sur $\operatorname{Spec} k$ : un torseur
    sous un $k$-module localement libre $V$ de rang $n+1$ ;
  \item[b)] d'applications $\varphi_i$ de $D_i$ dans l'ensemble des
    sous-espaces affines de dimension relative $i$ de $E$ ($0 \leq i \leq n$),
    qui respectent l'incidence ;
\end{enumerate}
telle que pour tout $g \in \operatorname{Aut}(\Pi)$ il existe un unique
automorphisme affine $\bar g$ de $E$ avec $\bar g \circ \varphi_i = \varphi_i
\circ g$ pour tout $i$ --- « cette unicité reste vraie fibre par
fibre »\footnote{Un troisième alinéa, « c) d'un homomorphisme de groupes
$\operatorname{Aut}(\Pi) \to \mathrm{Aut\,aff}(E)$ compatible avec\ldots »,
est biffé et remplacé par cette condition. La fin de la phrase qui annonce la
définition (« sur un anneau \ldots\ comme le \ldots ») n'est pas lue.}.

Il revient au même (page 6, « vérifier ») de se donner une application
$\Gamma$-équivariante $\varphi_R : R \to \operatorname{Drap}_n(E)$ dans les
drapeaux affines maximaux, et un homomorphisme $\varphi_\Gamma : \Gamma \to
\operatorname{Aff}(E)$ compatible, les sommets des drapeaux images engendrant
affinement $E$ --- ce qui rend $\varphi_\Gamma$ unique. Un repère origine $r$
étant fixé, il suffit de se donner $\varphi_\Gamma$ et le drapeau
$\varphi_R(r) = (f_0 \subset f_1 \subset \cdots \subset f_n)$, avec pour
seules conditions
\begin{enumerate}
  \item[a)] $\sigma_j f_i = f_i$ si $j \neq i$ ;
  \item[b)] les transformés $g f_0$, $g \in \Gamma$, engendrent affinement $E$
    fibre par fibre.
\end{enumerate}
En effet, a) dit que le stabilisateur $\mathfrak{G}_{\{i\}}$ de la facette
$i$-dimensionnelle de $r$ fixe $f_i$, de sorte que $\varphi_i(g\,\mathfrak{G}_{\{i\}})
= g f_i$ est bien défini, et b) donne l'unicité\footnote{La page identifie
$\operatorname{Aut}(\Pi)$ à $\Gamma$ (« $\mathfrak{G} = \Gamma$ ») ; cette
identification dépend du repère choisi et renverse l'ordre des produits, comme
le remarque la lecture du dossier 86 (page 14) et comme on le verra page 117.
La marge gauche de la page 6 porte une addition de quatre lignes, non lue.}.

\pagerange{7}{9}
\subsection*{Le repère des sommets (pages 7 à 9)}

\textbf{Proposition} (pages 7 et 8). \emph{Soit $s_0 = f_0$ et $s_{i+1} =
\sigma_i s_i$ ($0 \leq i \leq n$). Alors $(s_0, \dots, s_{n+1})$ est un repère
affine de $E$ fibre par fibre, $f_i$ est l'enveloppe affine de $s_0, \dots,
s_i$, et $\sigma_k s_i = s_i$ pour $0 \leq i < k \leq n$.}

On raisonne par récurrence sur $i$. Supposons $s_0, \dots, s_i$ affinement
libres, $f_i = \operatorname{Env}(s_0, \dots, s_i)$ et $\sigma_k s_j = s_j$ pour
$j \leq i < k$ --- c'est vrai pour $i = 0$, puisque les $\sigma_k$, $k \geq 1$,
fixent $f_0$. Comme $s_i \in f_i \subset f_{i+1}$ et que $\sigma_i$ fixe
$f_{i+1}$, on a $s_{i+1} = \sigma_i s_i \in f_{i+1}$ ; il suffit de voir que
$s_{i+1} \notin f_i$. Sinon, $\sigma_i f_i = \operatorname{Env}(\sigma_i f_{i-1},
\sigma_i s_i) = \operatorname{Env}(f_{i-1}, s_{i+1}) \subset f_i$, donc $\sigma_i
f_i = f_i$ ; comme les autres $\sigma_j$ fixent $f_i$ par a), $f_i$ serait
stable par $\Gamma$ et contiendrait l'orbite de $s_0$, contre b) (on a $f_i
\neq E$ pour $i \leq n$). Enfin, pour $k > i + 1$, $\sigma_k$ commute à
$\sigma_i$ et $\sigma_k s_{i+1} = \sigma_i \sigma_k s_i = \sigma_i s_i =
s_{i+1}$.

Il en découle $\sigma_i s_i = s_{i+1}$, $\sigma_i s_{i+1} = s_i$ (car
$\sigma_i^2 = 1$), et, pour $j \geq i + 2$, $\sigma_i s_j \in f_j \smallsetminus
f_{j-1}$ : $\sigma_i$ fixe $f_j$ et $f_{j-1}$, et $s_j \notin f_{j-1}$. Les
relations $\sigma_j f_i = f_i$ pour $j > i$ en résultent formellement ; pour
$j < i$ elles restent une condition, que la page note sans la
traiter\footnote{Dans le tableau de la page 9, $\sigma_1 s_4$ est dit dans $f_4
\smallsetminus f_2$ ; c'est $f_4 \smallsetminus f_3$, comme le schéma des autres
lignes. L'exposant de $\lambda_3^{0\,3}$ est lu avec doute.}.

\pagerange{9}{13}
\subsection*{Les paramètres (pages 9 à 13)}

Par les relations de commutation, les $\sigma_i$ sont déterminés par les
« valeurs critiques »
\[
\sigma_i s_{i+2} = s_0 + \lambda_{i1} e_1 + \cdots + \lambda_{i,i+2}\, e_{i+2}
\qquad (0 \leq i \leq n-1) ,
\]
qui forment une matrice à $n$ lignes et $n + 1$ colonnes, triangulaire, de
$2 + 3 + \cdots + (n+1) = n(n+3)/2$ coefficients. En effet, pour $j \geq i + 3$,
$s_j = \sigma_{j-1} \cdots \sigma_{i+2}\, s_{i+2}$, les $\sigma_k$ avec $k \geq
i + 2$ commutent à $\sigma_i$ et fixent $s_0, \dots, s_{i+1}$, d'où
\[
\sigma_i s_j = s_0 + \lambda_{i1} e_1 + \cdots + \lambda_{i,i+1} e_{i+1} +
\lambda_{i,i+2}\, e_j \qquad (j \geq i + 2) .
\]
Sur l'espace des translations $V$, cela donne, pour $1 \leq i \leq n-1$,
\[
\sigma_i e_j =
\begin{cases}
e_j & j < i \\
e_{i+1} & j = i \\
e_i & j = i + 1 \\
\lambda_{i1} e_1 + \cdots + \lambda_{i,i+1} e_{i+1} + \lambda_{i,i+2}\, e_j & j \geq i + 2 ,
\end{cases}
\]
\[
\sigma_0 e_j =
\begin{cases}
-e_1 & j = 1 \\
(\lambda_{01} - 1) e_1 + \lambda_{02}\, e_j & j \geq 2 ,
\end{cases}
\]
et $\sigma_n$ échange $s_n$ et $s_{n+1}$ en fixant les autres sommets. Les
déterminants sont $\det(\sigma_i|_V) = -(\lambda_{i,i+2})^{n-i}$ et
$\det(\sigma_0|_V) = -(\lambda_{02})^n$. Les conditions $\sigma_i^2 = 1$
s'écrivent
\[
\sigma_0^2 = 1 \iff (1 - \lambda_{01})(1 - \lambda_{02}) = 0,\ \lambda_{02}^2 = 1 ;
\]
\[
\sigma_i^2 = 1 \iff
\lambda_{ik}(1 + \lambda_{i,i+2}) = 0\ (k < i),\quad
\lambda_{i,i+1} + \lambda_{i,i+2}\lambda_{ii} = \lambda_{ii} + \lambda_{i,i+2}
\lambda_{i,i+1} = 0,\quad \lambda_{i,i+2}^2 = 1 ,
\]
et $\sigma_n^2 = 1$ toujours\footnote{La page 11 encadre déjà une forme réduite,
« $\sigma_i s_j = s_0 + \lambda_i(e_i + e_{i+1}) + e_j$ ». Le signe est faux :
avec $\lambda_{i,i+2} = 1$, l'involution exige $\lambda_{ii} = -\lambda_{i,i+1}$,
et la forme est $s_0 + \lambda_i(e_{i+1} - e_i) + e_j$, comme aux pages 18 et
25.}. Ce sont des calculs directs, qu'on a refaits.

\pagerange{13}{18}
\subsection*{Exiger des réflexions (pages 13 à 18)}

Ici intervient une exigence qui n'est pas dans la définition : chaque
$\sigma_i$ doit être une \emph{réflexion}\footnote{« Exprimons que $\sigma_0$ est
une réflexion », page 13 ; « Exprimons que $\sigma_i$ est une réflexion »,
page 15. Les conditions $\sigma_i^2 = 1$ seules laissent d'autres solutions :
$\lambda_{i,i+2} = -1$ donne, en caractéristique $\neq 2$, une involution dont
les points fixes forment un sous-espace de codimension $n - i + 1$. Pour un
polyèdre convexe réel, les générateurs sont des réflexions ; c'est sans doute
ce qui guide le choix.}. En coordonnées,
\[
\sigma_0(x) = \bigl(1 - x_1 + (\lambda_{01} - 1)\Sigma_0,\ \lambda_{02} x_2,\
\dots,\ \lambda_{02} x_{n+1}\bigr),
\]
dont les points fixes sont $2x_1 + (1 - \lambda_{01})\Sigma_0 = 1$,
$(\lambda_{02} - 1)x_j = 0$ ($j \geq 2$). Si $\lambda_{02} - 1$ est inversible,
c'est un point (ou le vide en caractéristique $2$), non un hyperplan ; sur un
corps, $\lambda_{02}^2 = 1$ force donc $\lambda_{02} = 1$. On pose
\[
\lambda_{02} = 1, \qquad \lambda_{01} = 1 + \alpha_0 ,
\]
et $\alpha_0$ est « l'invariant modulaire du polygone régulier induit par $\Pi$
dans $f_2$ ». De même, pour $1 \leq i \leq n-1$,
\[
\sigma_i(x) = \bigl(x_1 + \lambda_{i1}\Sigma_i,\ \dots,\ x_{i-1} +
\lambda_{i,i-1}\Sigma_i,\ x_{i+1} + \lambda_{ii}\Sigma_i,\ x_i +
\lambda_{i,i+1}\Sigma_i,\ \lambda_{i,i+2} x_{i+2},\ \dots,\ \lambda_{i,i+2}
x_{n+1}\bigr) .
\]
Si $\lambda_{i,i+2} = -1 \neq 1$, les points fixes vérifient $x_{i+2} = \cdots =
x_{n+1} = 0$ et $x_i = x_{i+1}$ : codimension $n - i + 1 \geq 2$. Donc
$\lambda_{i,i+2} = 1$. Si alors l'un des $\lambda_{ik}$, $k < i$, n'est pas nul,
les points fixes vérifient $\Sigma_i = 0$ et $x_i = x_{i+1}$ : codimension $2$.
Donc $\lambda_{ik} = 0$ pour $k < i$, et il reste $\lambda_{ii}\Sigma_i = x_i -
x_{i+1} = -\lambda_{i,i+1}\Sigma_i$, qui est un hyperplan si et seulement si
$\lambda_{ii} + \lambda_{i,i+1} = 0$ --- c'est aussi ce que demande $\sigma_i^2
= 1$. On pose $\lambda_{i,i+1} = 1 + \alpha_i$. Sur un anneau quelconque, la
même conclusion vaut dès qu'on demande que le lieu fixe soit un hyperplan
relatif, défini par une seule équation de partie linéaire unimodulaire : chaque
autre équation doit en être multiple, et la comparaison des coefficients de
$x_1$ (pour $\sigma_0$) ou de $x_i$ (pour $\sigma_i$) donne $\lambda_{02} = 1$,
$\lambda_{ik} = 0$, $\lambda_{i,i+2} = 1$ et $\lambda_{ii} = -\lambda_{i,i+1}$
\footnote{La page traite les anneaux autrement : « ces résultats marchent
aussi sur un anneau réduit » ; sur un anneau quelconque elle pose
$\lambda_{i,i+2} = 1 + \nu$, $\nu$ nilpotent, distingue la caractéristique
résiduelle $\neq 2$ et $= 2$ (système II, page 19), et conclut « manifestement »
pour le second cas. L'argument par comparaison des coefficients est de
l'édition. Une ligne de la page 18 porte $\lambda_{i,i+1} + \lambda_{i,i+1}$,
pour $\lambda_{ii} + \lambda_{i,i+1}$. Pour $\sigma_0$, si la caractéristique
est $2$ et $\alpha_0 = 0$, l'équation $2x_1 - \alpha_0\Sigma_0 = 1$ est vide :
$\sigma_0$ est la translation de vecteur $e_1$, et son hyperplan fixe est
l'hyperplan à l'infini (« l'hyperplan : l'infini »).}.

\textbf{Forme normale} (pages 15 et 18). \emph{Une réalisation affine régulière
de $\Pi$ dont les générateurs sont des réflexions s'écrit, dans le repère des
sommets,}
\[
\begin{aligned}
\sigma_0(x) &= \bigl(1 - x_1 + \alpha_0\Sigma_0,\ x_2,\ \dots,\ x_{n+1}\bigr), \\
\sigma_i(x) &= \bigl(x_1, \dots, x_{i-1},\ x_{i+1} - (1 + \alpha_i)\Sigma_i,\
x_i + (1 + \alpha_i)\Sigma_i,\ x_{i+2}, \dots, x_{n+1}\bigr) \quad (1 \leq i
\leq n-1), \\
\sigma_n(x) &= \bigl(x_1, \dots, x_{n-1},\ x_{n+1},\ x_n\bigr),
\end{aligned}
\]
\emph{de miroirs}
\[
H_0 : 2x_1 - \alpha_0\Sigma_0 = 1, \qquad
H_i : x_i - x_{i+1} + (1 + \alpha_i)\Sigma_i = 0, \qquad
H_n : x_n = x_{n+1} .
\]
\emph{Elle dépend de $n$ paramètres $\alpha_0, \dots, \alpha_{n-1}$.}\footnote{La
page écrit l'équation de $H_i$ avec $x_{n+2}$, pour $x_{n+1}$. Le signe devant
$\alpha_0$ dans $H_0$, page 15, est lu avec doute ; le calcul donne $-$.}

\pagerange{19}{22}
\subsection*{Les rotations élémentaires, et une question (pages 19 à 22)}

\emph{Relations.} Inversement, les $\sigma_i$ de la forme normale vérifient les
relations de $\mathfrak{G}_n$ quels que soient les $\alpha_i$. Ce sont des
involutions, et la page 19 observe que $\sigma_0$ ne change que la première
coordonnée et $\sigma_i$ que la $i$-ième et la $(i+1)$-ième, ce qui
« implique » que $\sigma_i$ et $\sigma_j$ commutent pour $j \geq i + 2$. Il faut
y ajouter que $\sigma_j$ conserve la somme $x_j + x_{j+1}$, donc $\Sigma_i$ et
$\Sigma_0$, dont dépend $\sigma_i$\footnote{L'ajout est de l'édition ; on a
vérifié les relations sur machine pour $n \leq 4$.}.

\emph{Les rotations.} La page calcule $\sigma_0\sigma_1(x) = \bigl(1 + \alpha_0
x_1 - x_2 + (1 + \alpha_1 + 2\alpha_0 + \alpha_0\alpha_1)\Sigma_1,\ x_1 + (1 +
\alpha_1)\Sigma_1,\ x_3, \dots\bigr)$, écrit $(\sigma_0\sigma_1)^p = (F_p, G_p,
x_3, \dots)$ et s'arrête sur « Donc $(\sigma_0\sigma_1)^p = \mathrm{id} \iff$ ».
La page 20 reprend la question par les quotients successifs du drapeau :
$\sigma_i$ est l'identité sur $f_{i-1}$ et sur $E / f_{i+1}$, et échange deux
coordonnées sur $f_{i+1} / f_{i-1}$ ; pour $\sigma_i\sigma_{i+1}$ sur $f_{i+2} /
f_{i-1}$ et pour $\sigma_0\sigma_1$ sur $f_2$, la page laisse le membre de droite
en blanc. La page 22 calcule $\pi = \sigma_0\sigma_1$, $\pi^2$, $\pi^3$ sur
$f_2$ et trouve, pour un triangle, $\alpha_0 = -1$. Voici la réponse
complète, qui est de l'édition.

\textbf{Proposition.} \emph{La partie linéaire de $\sigma_{i-1}\sigma_i$ ($1 \leq i
\leq n$) a pour polynôme caractéristique}
\[
(t - 1)^{n-1}\,(t^2 - \alpha_{i-1} t + 1) .
\]
\emph{Sur un corps $k$, si $3 \leq p_i < \infty$ et si $p_i$ est inversible dans
$k$, $\sigma_{i-1}\sigma_i$ est d'ordre exactement $p_i$ si et seulement si
$\alpha_{i-1} = \zeta + \zeta^{-1}$ avec $\zeta$ racine primitive $p_i$-ième de
l'unité (dans $k$ ou dans une extension quadratique), c'est-à-dire si
$\alpha_{i-1}$ est racine de $\Psi_{p_i}$.}

Sur $V(f_{i+1}) / V(f_{i-2})$, de base les images de $e_{i-1}, e_i, e_{i+1}$,
$\sigma_{i-1}\sigma_i$ envoie $e_{i-1} \mapsto e_i$, $e_i \mapsto -(1 +
\alpha)e_{i-1} + (1 + \alpha)e_i + e_{i+1}$, $e_{i+1} \mapsto e_{i-1}$ ($\alpha =
\alpha_{i-1}$), de polynôme caractéristique $(t - 1)(t^2 - \alpha t + 1)$ ; il est
l'identité sur $V(f_{i-2})$ et sur $V / V(f_{i+1})$ (pour $i = 1$ : la matrice
$\bigl(\begin{smallmatrix} \alpha_0 & -1 \\ 1 & 0 \end{smallmatrix}\bigr)$ sur
$\langle e_1, e_2 \rangle$). Comme $\sigma_{i-1}\sigma_i - 1$ est de rang $\leq
2$, la valeur propre $1$ a un sous-espace propre de dimension $\geq n - 1$ ; si
$\alpha \neq \pm 2$, les racines $\zeta^{\pm 1}$ sont distinctes et $\neq 1$,
la partie linéaire est diagonalisable, et son ordre est celui de $\zeta$. Pour
$i \geq 2$, $\sigma_{i-1}\sigma_i$ fixe $s_0$ ; pour $i = 1$, il a un point fixe
dans $f_2$, parce que $1$ n'est pas valeur propre sur $\langle e_1, e_2 \rangle$.
Si l'ordre est $p_i$, premier à la caractéristique, la partie linéaire est
semi-simple et ses valeurs propres sont des racines $p_i$-ièmes de l'unité ;
$\zeta = \pm 1$ donnerait un ordre $\leq 2$\footnote{Machine : le polynôme
caractéristique pour $n = 2, 3, 4$ et tout $i$. L'énoncé pour $n = 2$ est celui
de la page 67, « $\rho_s^p = \mathrm{id} \iff F_p(\alpha) = 0$ », sans
démonstration.}.

Ainsi, \emph{à choix du repère près, les réalisations affines régulières par
réflexions du polytope $\{p_1, \dots, p_n\}$ sur un corps où les $p_i$ sont
inversibles correspondent aux suites $(\alpha_0, \dots, \alpha_{n-1})$ telles
que $\Psi_{p_i}(\alpha_{i-1}) = 0$.} Sur les réels, le choix $\alpha_{i-1} =
2\cos(2\pi/p_i)$ donne le polytope convexe ; les autres racines, $2\cos(2\pi
k/p_i)$, donnent les polytopes étoilés (page 69). Pour $p_i = \infty$, la page 71
demande $\alpha_{i-1} = 2$ : la rotation est alors unipotente. La partie
linéaire est une représentation du groupe de Coxeter engendrée par des
réflexions ; pour un groupe fini, c'est, à un changement de base près, la
représentation géométrique de Tits ou l'une de ses conjuguées
galoisiennes\footnote{Le rapprochement est le nôtre ; dans la représentation de
Tits, la forme invariante a pour coefficients $-\cos(\pi / m_{ij})$, et l'on a
$\alpha + 2 = 4\cos^2(\pi / p)$ pour $\alpha = 2\cos(2\pi / p)$.}.

\textbf{Une question} (page 21). Supposons $k$ un corps et soit $G$ l'image de
$\Gamma$ dans $\operatorname{Aut}_{\mathrm{aff}}(E)$. La représentation de $G$ est fidèle,
et $\varphi_R : R \to \operatorname{Drap}(E)$ injective. Les applications $D_i(\Pi)
\to \operatorname{Grass}_i^{\mathrm{aff}}(E)$ le sont-elles aussi ? Autrement
dit, si $g, g' \in G$ vérifient $g f_i = g' f_i$, $g^{-1} g'$ est-il dans le
sous-groupe engendré par les $\sigma_j$, $j \neq i$ ? « Ce n'est pas évident du
tout », et la page soupçonne un lien avec des résultats sur les longueurs dans
la réalisation riemannienne\footnote{La fin de la phrase est en partie
incertaine (« des \emph{Longueurs} \emph{associés} à la réal.\
\emph{riemannienne} »). Deux lignes biffées ne se lisent que par morceaux. Le
$\Pi$ de la page est une « $n$-carte », avec un mot incertain.}. La question
n'est pas tranchée ; pour l'icosaèdre, la réponse est oui en toute
caractéristique (pages 93 et 94, et ci-dessous).

\pagerange{23}{25}
\subsection*{La forme normale, en coordonnées barycentriques (pages 23 à 25)}

Les pages 23 à 25 récrivent la forme normale de façon homogène. Soit $\ell_i$
la forme affine
\[
\ell_0 = 1 - 2x_1 + \alpha_0\Sigma_0, \qquad
\ell_i = x_i - x_{i+1} + (1 + \alpha_i)\Sigma_i \quad (1 \leq i \leq n-1),
\qquad \ell_n = x_n - x_{n+1} ,
\]
c'est-à-dire, en coordonnées barycentriques $(t_0, \dots, t_{n+1})$ relatives à
$(s_0, \dots, s_{n+1})$, $\ell_i = t_i - t_{i+1} + (1 + \alpha_i)(t_{i+2} +
\cdots + t_{n+1})$ pour tout $i$ (la somme est vide pour $i = n$). Alors
\[
\sigma_i(P) = P + \ell_i(P)\,(s_{i+1} - s_i) \qquad (0 \leq i \leq n) ,
\]
et, comme $\ell_i(s_{i+1}) - \ell_i(s_i) = -2$, $\sigma_i$ est une involution
de miroir $H_i = \{\ell_i = 0\}$ et de « centre » $\delta_i$, le point à
l'infini de la direction $s_{i+1} - s_i$ : une réflexion projective, au sens
de la géométrie projective classique. La page écrit
$\sigma_i - \mathrm{id} = \ell_i \otimes (s_{i+1} - s_i)$ en termes de la base
duale $(s'_j)$, et, en homogénéisant par une variable $z$, $\bar\sigma_0(\bar
x) = \bar x + [e'_0 - 2e'_1 + \alpha_0(e'_2 + \cdots + e'_{n+1})] \otimes
e_1$\footnote{La page 25 écrit $\sigma_0 s_i = -(1 + \alpha_i)s_0 + (1 +
\alpha_i)s_1 + s_i$ ; l'indice de $\alpha$ est $0$, comme le donne le tableau
encadré qui suit, valable pour $0 \leq i \leq n$ : $\sigma_i s_j = -(1 +
\alpha_i)s_i + (1 + \alpha_i)s_{i+1} + s_j$ pour $j \geq i + 2$. Le coefficient
$2$ de $e'_1$, page 23, est une surcharge lue d'après $H_0$.}.
La page 23 commence aussi d'écrire une forme quadratique générale en $x_1,
\dots, x_{n+1}$, sous le titre « Invariants ? », et la biffe : la question est
reprise, pour $n = 2$, à la page 67.

\pagerange{28}{36}
\section*{II. Les $n$-quasi-polyèdres (pages 28 à 36)}

La chemise de la page 28 porte « $n$-quasi-polyèdres (vieilles rédactions) ».
On abandonne les relations $\sigma_i^2 = 1$.

\pagerange{29}{29}
\subsection*{Le repère, sans involutions (page 29)}

Soient $\sigma_0, \dots, \sigma_n \in \operatorname{Aff}(E)$, $E$ affine de rang
$n+1$, avec $\sigma_i\sigma_j = \sigma_j\sigma_i$ pour $j \geq i + 2$, et un
drapeau $E_0 = \{s_0\} \subset E_1 \subset \cdots \subset E_n$ tel que $\sigma_i
E_j = E_j$ pour $j \neq i$ et $\sigma_i E_i \neq E_i$. On pose $s_{i+1} =
\sigma_i s_i$.

\textbf{Proposition} (page 29). \emph{$(s_0, \dots, s_{n+1})$ est un repère
affine de $E$, $E_i = \operatorname{Env}(s_0, \dots, s_i)$ et $\sigma_k s_i = s_i$
pour $0 \leq i < k \leq n$.}

La démonstration est celle de la page 7, l'hypothèse $\sigma_{i-1}E_{i-1} \neq
E_{i-1}$ remplaçant la génération par l'orbite de $s_0$ : si $s_i =
\sigma_{i-1}s_{i-1}$ était dans $E_{i-1}$, on aurait $\sigma_{i-1}E_{i-1} \subset
E_{i-1}$, donc l'égalité\footnote{La page lit l'enveloppe affine « Env »,
sans certitude sur les dernières lettres. Ici $E$ est sur un corps, ou fibre
par fibre.}.

\pagerange{30}{33}
\subsection*{Les $2n + 1$ paramètres (pages 30 à 33)}

Il suffit de connaître $u_i = \sigma_i s_{i+1} = \sigma_i^2 s_i \in E_{i+1}$ et
$v_i = \sigma_i s_{i+2} \in E_{i+2}$, puisque $\sigma_i s_j = \sigma_{j-1}\cdots
\sigma_{i+2}(v_i)$ pour $j > i + 2$. Avec $u_i = s_0 + \sum_{k \leq i+1}
u_{ik}e_k$ et $v_i = s_0 + \sum_{k \leq i+2} v_{ik}e_k$,
\[
\sigma_i s_j = s_0 + \sum_{k \leq i+1} v_{ik}e_k + v_{i,i+2}\,e_j \quad (j \geq i
+ 2),
\]
et l'on écrit les matrices. Pour $\sigma_0$ :
\[
\sigma_0(x) = \bigl(1 + (u_{01} - 1)x_1 + (v_{01} - 1)\Sigma_0,\ v_{02}x_2,\
\dots,\ v_{02}x_{n+1}\bigr),
\]
dont le lieu fixe est un hyperplan (projectif) si et seulement si $v_{02} =
1$\footnote{Sans l'involution, « réflexion projective » veut dire ici : une
transformation qui fixe point par point un hyperplan --- une homologie ou une
élation, dirait-on en géométrie projective classique. Le mot de la page est
« réflexion projective ».} ; on pose $u_{01} = \lambda_0$, $v_{01} = \mu_0$.
Pour $1 \leq i \leq n - 1$,
\[
\sigma_i(x) = \bigl(x_1 + u_{i1}x_{i+1} + v_{i1}\Sigma_i,\ \dots,\ u_{ii}x_{i+1} +
v_{ii}\Sigma_i,\ x_i + u_{i,i+1}x_{i+1} + v_{i,i+1}\Sigma_i,\
v_{i,i+2}x_{i+2},\ \dots\bigr),
\]
et le lieu fixe est un hyperplan si et seulement si les autres équations sont
conséquences de la $i$-ième, soit
\[
u_{ii} + u_{i,i+1} = 1,\quad v_{ii} + v_{i,i+1} = 0,\quad v_{i,i+2} = 1,\quad
u_{ik} = v_{ik} = 0\ (k < i) .
\]
Avec $\lambda_i = u_{ii}$, $\mu_i = v_{ii}$ :
\[
\sigma_i(x) = \bigl(x_1, \dots, x_{i-1},\ \lambda_i x_{i+1} + \mu_i\Sigma_i,\
x_i + (1 - \lambda_i)x_{i+1} - \mu_i\Sigma_i,\ x_{i+2}, \dots, x_{n+1}\bigr) ;
\]
$\sigma_n$ n'a que $\lambda_n$ ($\Sigma_n = 0$). Le cas général a donc $2n + 1$
paramètres, $\lambda_0, \dots, \lambda_n$ et $\mu_0, \dots, \mu_{n-1}$. Enfin
\[
\sigma_i^2 = \mathrm{id} \iff \lambda_i = 1 \ (i \geq 1), \qquad
\sigma_0^2 = \mathrm{id} \iff \lambda_0 = 0 ,
\]
et l'on retrouve le cas polyédral, à $n$ paramètres $\mu_i$\footnote{Le calcul
de $\sigma_i^2$ de la page 33 omet le terme $(1 - \lambda_i)\mu_i\Sigma_i$ de
$X_i$ et donne à $x_{i+1}$, dans $X_{i+1}$, le coefficient $1$ au lieu de
$\lambda_i^2 - \lambda_i + 1$ ; la conclusion « sss $\lambda_i = 1$ » n'en est pas
affectée. Pour $\sigma_0$, $X_1 = \lambda_0 + (\lambda_0 - 1)^2x_1 +
\lambda_0(\mu_0 - 1)\Sigma_0$, exact. Dans le cas polyédral, $\mu_0 - 1 =
\alpha_0$ et $\mu_i = -(1 + \alpha_i)$ ($i \geq 1$) en termes de la page 18.}.

\pagerange{34}{36}
\subsection*{Les ordres (pages 34 à 36)}

Pour $n = 1$ (un « quasi-polygone » du plan), avec $\eta_0 = \lambda_0 - 1$,
$\alpha_0 = \mu_0 - 1$ :
\[
\sigma_0(x, y) = (1 + \eta_0 x + \alpha_0 y,\ y), \qquad \sigma_1(x, y) =
\bigl(\lambda_1 y,\ x + (1 - \lambda_1)y\bigr),
\]
sur un corps $k$ de caractéristique $p \geq 0$. La partie linéaire de
$\sigma_0$ a pour valeurs propres $1$ et $\eta_0$ ; ses puissances sont
$\bigl(\begin{smallmatrix} \eta_0^m & \alpha_0 S_m \\ 0 & 1 \end{smallmatrix}\bigr)$
avec $S_m = 1 + \eta_0 + \cdots + \eta_0^{m-1}$, et $\sigma_0^m(x, 0) = (S_m +
\eta_0^m x, 0)$. Donc
\[
\sigma_0^m = \mathrm{id} \iff S_m(\eta_0) = 0 ,
\]
et l'ordre de $\sigma_0$ est celui de $\eta_0$ dans $k^*$ si $\eta_0 \neq 1$. Si
$\eta_0 = 1$, $\sigma_0(x, y) = (x + 1 + \alpha_0 y, y)$, de partie linéaire la
transvection $\bigl(\begin{smallmatrix} 1 & \alpha_0 \\ 0 & 1
\end{smallmatrix}\bigr)$ : $\sigma_0$ est d'ordre $p$ si $p > 0$, infini si $p =
0$, quel que soit $\alpha_0$\footnote{La page distingue a) « $1 + \eta_0 +
\cdots + \eta_0^{n-1} = 0$, i.e.\ $\eta_0^n = 1$ et ($\eta_0 \neq 1$ ou $p \nmid
n$) » et b) « $\eta_0 = 1$, $\alpha_0 = 0$ \ldots\ et $p \mid n$ ». Les
divisibilités sont échangées : a) est « ($\eta_0 \neq 1$ et $\eta_0^n = 1$) ou
($\eta_0 = 1$ et $p \mid n$) », et le cas restant où la partie linéaire est
l'identité est b) $\eta_0 = 1$, $\alpha_0 = 0$, $p \nmid n$ ; la lecture de la
parenthèse de a) est d'ailleurs incertaine. Page 35, sous « $\eta_0 = 1$ », la
page écrit « a) $\alpha_0 = 1$ : $\sigma_0 = \mathrm{id}$, cas impropre » : c'est
impossible, $\sigma_0$ envoie $s_0$ sur $s_1$ ; la ligne précédente dit
« d'ordre $1$ sss $\alpha_0 = 0$ », ce qui concerne la partie linéaire.}.
De même $\sigma_1$ a pour valeurs propres $1$ et $\delta_1 = -\lambda_1$ ; son
ordre est celui de $\delta_1$ sauf si $\delta_1 = 1$, où $\sigma_1(x, y) = (-y, x
+ 2y)$ est unipotent non trivial, d'ordre infini si $p = 0$, $p$ sinon.

En dimension $n$ (page 36), avec $\eta_i = -\lambda_i$ ($i \geq 1$), la partie
linéaire de chaque $\sigma_i$ a la valeur propre $1$ avec multiplicité $n$ et
la valeur propre $\eta_i$ ; $\sigma_i - \mathrm{id}$ étant de rang $1$, elle est
diagonalisable si $\eta_i \neq 1$, et l'ordre de $\sigma_i$ est alors celui de
$\eta_i$, comme la page le demande avec un point
d'interrogation\footnote{La page écrit « $1$ de multiplicité $n - 1$ » ;
l'espace est de dimension $n + 1$. Elle récrit aussi $\sigma_i$ avec un
paramètre qu'elle nomme $\alpha_i$, $\mu_i = 1 + \alpha_i$, qui n'est pas celui
de la page 18 : dans le cas polyédral, le $1 + \alpha_i$ de la page 36 est le
$-(1 + \alpha_i)$ de la page 18. On n'emploie pas cette lettre ici.}.

\pagerange{38}{47}
\section*{III. « Étude locale » des $n$-polyèdres et des $n$-cartes (pages 38 à 47)}

La chemise de la page 38 porte « “étude locale” des $n$-polyèdres et des
$n$-cartes (vieilles rédactions) ».

\pagerange{39}{42}
\subsection*{Bord, étoile, sections (pages 39 à 42)}

Soit $C$ une $n$-carte. Pour $0 \leq i \leq n$, notons $\mathfrak{G}^{<i} =
\langle \sigma_0, \dots, \sigma_{i-1} \rangle$ et $\mathfrak{G}^{>i} = \langle
\sigma_{i+1}, \dots, \sigma_n \rangle$ ; ils commutent, et $D_i = C /
(\mathfrak{G}^{<i} \times \mathfrak{G}^{>i})$.
\begin{itemize}
  \item Sur $C / \mathfrak{G}^{>i} \to D_i$ opère $\mathfrak{G}^{<i} \simeq
    \mathfrak{G}_{i-1}$, dans les fibres : c'est une $(i-1)$-carte dont les
    composantes connexes correspondent aux $i$-facettes. La composante de $f_i$
    est le \emph{bord} $\partial f_i$ ; ses $j$-facettes sont les drapeaux de
    type $(j, i)$ de composante $f_i$, et plus généralement ses drapeaux de
    type $(j_1 < \cdots < j_p)$ sont les drapeaux de type $(j_1, \dots, j_p,
    i)$ de $C$ au-dessus de $f_i$.
  \item Dualement, sur $C / \mathfrak{G}^{<i} \to D_i$ opère $\mathfrak{G}^{>i}
    \simeq \mathfrak{G}_{n-i-1}$ : une $(n-i-1)$-carte, l'\emph{étoile}
    $\delta f_i$, dont les facettes de dimension $j$ sont les drapeaux de type
    $(i, i + j + 1)$ au-dessus de $f_i$\footnote{La page note $E^{-}(i)$,
    $E^{+}(i)$, $\Gamma_n^{\pm}(i)$ avec des signes repassés ; telle qu'elle
    est lue, elle quotiente par le groupe inférieur et fait opérer
    $\Gamma_n^{+}(i) \simeq \Gamma_{i-1}$, ce que la note de la transcription
    signale. On suit les définitions : on quotiente par le groupe supérieur, le
    groupe inférieur opère. La carte est appelée « $n$-quelette », mot lu avec
    doute.}.
  \item Plus généralement (pages 40 et 41), pour $0 \leq i$, $i + 2 \leq j \leq
    n + 1$ (avec les conventions $f_{-1}$, $f_{n+1}$), $C / (\mathfrak{G}^{<i}
    \times \mathfrak{G}^{>j}) = D_{\{i, i+1, \dots, j\}} \to D_{\{i, j\}}$, sur
    lequel opère $\langle \sigma_{i+1}, \dots, \sigma_{j-1} \rangle \simeq
    \mathfrak{G}_{j-i-2}$ : la \emph{section} $\Pi(d_{ij})$ au-dessus d'un
    drapeau $d_{ij} = (f_i < f_j)$, une $(j - i - 2)$-carte dont les facettes de
    dimension $h$ sont les drapeaux de type $(i, i + h + 1, j)$. Pour $j = i
    + 2$, c'est un ensemble à deux éléments sous $\sigma_{i+1}$ quand la carte
    est un polyèdre ; pour $j = i + 3$, un polygone, de nombre de côtés
    $\nu(d_{ij}) \in [1, +\infty]$\footnote{Le cas $j = i + 2$ porte deux mots
    lus « un pas idiot », lecture très incertaine ; on comprend « cas idiot ».
    Le bord et l'étoile sont les sections $\Pi(d_{-1,i})$ et $\Pi(d_{i,n+1})$,
    comme le dit la page 41.}.
\end{itemize}
Aujourd'hui, ce sont les \emph{sections} $F/G$ d'un polytope abstrait ; le bord
est la section $f_i/f_{-1}$, l'étoile la « figure de facette » $f_{n+1}/f_i$.

\textbf{Les invariants} (page 42). Ce sont le nombre de côtés des $2$-faces,
$\nu(\partial f_2)$, ceux des sections $\Pi(d_{k-1,k+2})$ ($1 \leq k \leq n-2$),
et celui de l'étoile des $(n-2)$-facettes. Pour $\Pi$ régulier, ce sont des
constantes, et la page affirme qu'elles valent les ordres des
$\sigma_k\sigma_{k+1}$ --- « vrai tout au moins dans le cas des géométries
d'incidence »\footnote{La page numérote ces invariants $p_0, \dots, p_{n-1}$ ;
dans la convention de la page 3, suivie ici, ce sont $p_1, \dots, p_n$. Son
début de démonstration écrit $(\sigma_0\sigma_1)^p x = g_x x$ avec $g_x \in
G_n^{+}(2)$.}. La réserve est juste : pour une carte régulière $C = \Gamma$,
le nombre de côtés de la section est l'ordre de $\sigma_k\sigma_{k+1}$ agissant
sur l'orbite correspondante, qui est l'ordre dans $\Gamma$ dès que $\langle
\sigma_k, \sigma_{k+1}\rangle$ rencontre trivialement le stabilisateur ; c'est
une conséquence de la propriété d'intersection, qui vaut pour les polytopes et
peut manquer pour une carte quelconque.

\pagerange{43}{45}
\subsection*{Représentation géométrique (pages 43 à 45)}

Soit $\Pi$ une $n$-carte, $R = \operatorname{Rep}(\Pi)$ son $\mathfrak{G}_n$-ensemble
de repères. Une \emph{représentation géométrique} de $\Pi$ dans un espace affine
$E$ sur un corps est une application
\[
\varphi_R : R \to \mathrm{Drap\,max}(E)
\]
dans l'ensemble des drapeaux $(E_0 \subset E_1 \subset \cdots \subset E_n
\subset E)$, $E_i$ de rang $i$, telle que $\pi_j \varphi_R(\sigma_i r) = \pi_j
\varphi_R(r)$ pour $i \neq j$, où $\pi_j$ prend la composante de dimension $j$.
Par passage au quotient, $\pi_j\varphi_R$ se factorise en $\varphi_j : D_j(\Pi) \to
\operatorname{Grass}^{\mathrm{aff}}_j(E)$\footnote{Le corps est dit « ordonné »,
lecture incertaine. La page 44, barrée, est un premier essai, avec un
diagramme $\varphi_G : G \to \operatorname{Aff}(E)$, $\varphi_R : R \to
\mathrm{Drap\,max}(E)$, $\varphi_{D_i} : D_i(\Pi) \to
\operatorname{Grass}_i(E)$, où $R$ est un « $G$-torseur » et $\varphi_R$ un
« $G$-homomorphisme ». La version de la page 43 n'exige pas de groupe : c'est
une réalisation, régulière ou non.}.

Lorsque $\Pi$ est un polyèdre (une géométrie d'incidence), se donner $\varphi_R$
revient à se donner les $\varphi_j$, compatibles à l'incidence. En général, il
faut tous les
\[
\varphi_J : D_J(\Pi) \to \mathrm{Drap\,aff}_J(E), \qquad J \subset \Delta_n
= \{0, \dots, n\},
\]
compatibles aux restrictions pour $J' \subset J$ : un morphisme de « géométries
de drapeaux » de base $\Delta_n$, de $D(\Pi)$ dans $\mathrm{Drap\,aff}(E)$
--- et c'est la donnée de $\varphi_R$\footnote{La page écrit le but
$\mathrm{Drap\,aff}_J(\Pi)$ ; on lit $(E)$. Les géométries de drapeaux
sont l'objet des pages 15 à 33 du dossier 76.}. En effet, $R = D_{\Delta_n}$,
et $D_J$ est le quotient de $R$ par $\mathfrak{G}_J$.

\pagerange{46}{47}
\subsection*{Réalisations induites (pages 46 et 47)}

Soit $f_i \in D_i(\Pi)$, $1 \leq i \leq n$, et $E_i = \varphi_i(f_i)$. Le bord
$\partial f_i$ a pour ensemble de repères l'image réciproque de $\{f_i\}$ par
$D_{\{0, \dots, i\}}(\Pi) \to D_i(\Pi)$, et les données de $\varphi$ en
définissent une réalisation canonique, dite induite, dans $E_i$. « Légèrement
moins évident » est le cas de l'étoile $\delta f_i$ : ses sommets sont les
drapeaux $(f_i < f_{i+1})$, qui s'envoient sur des couples $E_i \subset
E_{i+1}$, et si $V_i$ est l'espace des translations de $E_i$, ces $E_{i+1}$
correspondent aux droites de $E / V_i$ passant par le point image de $E_i$ ---
c'est-à-dire aux points de l'espace projectif $\mathbb{P}(V / V_i)$. La page
s'arrête là\footnote{La première moitié de la page 46 est barrée ; elle
définissait la carte quotient $R_{ij} = R / (\sigma_0, \dots, \sigma_{i-1},
\sigma_{j+1}, \dots, \sigma_n)$, avec $\sigma'_0 = \sigma_{i+1}, \dots,
\sigma'_{j-i-2} = \sigma_{j-1}$. Le mot de liaison en tête de la page 47 n'est
pas lu ; la phrase finale sur $\mathbb{P}(V / V_i)$ est une conclusion de
l'édition.} : l'étoile se réalise naturellement non dans un espace affine, mais
dans un espace projectif --- comme la figure de sommet d'un polyèdre convexe,
vue depuis le sommet.

\pagerange{49}{59}
\section*{IV. Les $2$-polyèdres repérés (pages 49 à 59)}

Deux chemises : « 2 -- Polyèdres réguliers. (Vieille rédaction) » (page 49), et
« (18) 2-Polyèdres repérés en coordonnées universelles (vieille rédaction) »
(page 50).

\pagerange{51}{53}
\subsection*{Réalisation géométrique d'un polyèdre combinatoire (pages 51 à 53)}

Soit $X = (\mathrm{S}, \mathrm{A}, \mathrm{F}, \mathrm{R}, \mathrm{R}',
\mathrm{R}'')$ un polyèdre combinatoire, $\mathrm{R}$ l'ensemble de ses repères,
et $E$ un espace affine (ou projectif) sur un corps $k$. Une
\emph{réalisation géométrique} de $X$ est la donnée de $\varphi_S : \mathrm{S} \to
E$, $\varphi_A : \mathrm{A} \to \operatorname{Dr}(E)$ (droites), $\varphi_F :
\mathrm{F} \to \operatorname{Pl}(E)$ (plans), compatibles aux incidences, telles
que
\begin{enumerate}
  \item[a)] les deux sommets d'une arête ont des images distinctes, et de même
    les deux faces qui la bordent ;
  \item[b)] si $s < f$, les deux arêtes entre $s$ et $f$ ont des images
    distinctes.
\end{enumerate}
Alors l'une quelconque des trois applications détermine les deux autres.

\textbf{Proposition} (pages 51 et 52). \emph{Soit $r_0 = (s_0, a_0, f_0)$ un
repère, $s_1$ le second sommet de $a_0$, $a_1$ la seconde arête entre $s_0$ et
$f_0$, $s_2$ son second sommet, $f_1$ la seconde face qui borde $a_0$, $a_2$ la
seconde arête entre $s_0$ et $f_1$, et $s_3$ son second sommet. Les images de
$s_0, s_1, s_2, s_3$ sont affinement indépendantes.}

$s_1 \neq s_0$ et $s_2 \neq s_0$ par a) ; $a_1 \neq a_0$ par b), donc $s_2$
n'est pas sur la droite $a_0$, et $f_0$ est le plan $(s_0, s_1, s_2)$. De même
$f_1$ est le plan $(s_0, s_1, s_3)$ ; comme $f_1 \neq f_0$ par a), $s_3 \notin
f_0$. D'où (corollaire) : $a_0 = \operatorname{dr}(s_0, s_1)$, $f_0 =
\operatorname{pl}(s_0, s_1, s_2)$, et $\varphi_S$ détermine $\varphi_A$,
$\varphi_F$ ; dualement $a_0 = f_0 \cap f_1$, $s_0 = f_0 \cap f_1 \cap f_2$ ($f_2$
la seconde face qui borde $a_1$), et $\varphi_F$ détermine les autres ; enfin
$s_0 = a_0 \cap a_1$, $f_0 = \operatorname{pl}(a_0, a_1)$. Si $E$ est de
dimension $3$, $\varphi_S(\mathrm{S})$ engendre $E$\footnote{La page conclut
« $\varphi_S(\mathrm{S})$ engendre aff.\ $E$ » sans dire que $E$ est de
dimension $3$ ; on l'ajoute. L'énoncé dual (trois plans indépendants) est
projectif. La page 59, sans texte, redessine la figure de la page 52 :
le sommet $s_0$, les arêtes $a_0, a_1, a_2$ vers $s_1, s_2, s_3$, les faces
$f_0, \dots, f_3$.}.

\textbf{Réalisation régulière.} $X$ étant régulier, de groupe $\mathrm{G} =
\operatorname{Aut}(X)$, la réalisation est \emph{régulière} si pour tout
$\mathrm{g} \in \mathrm{G}$ il existe $g \in \operatorname{Aut}(E)$ (affine ou
projectif) avec $g \circ \varphi_S = \varphi_S \circ \mathrm{g}$, et de même pour
$\varphi_A$, $\varphi_F$ --- l'une de ces relations entraînant les autres. Dans le
cas affine, $g$ est unique et $\mathrm{g} \mapsto g$ est un
homomorphisme\footnote{Ou un anti-homomorphisme, selon qu'on fait opérer
$\mathrm{G}$ à gauche ou à droite sur les repères ; voir la page 117.}. Dans le
cas projectif, c'est encore vrai pourvu qu'il existe un sommet dont l'image
n'est dans aucun des quatre plans engendrés par trois des points $s_0, s_1,
s_2, s_3$, car alors les images contiennent un repère projectif ; l'exception
est celle où $\mathrm{S} = \{s_0, s_1, s_2, s_3\}$, le tétraèdre, dont les
automorphismes projectifs fixant les quatre sommets forment un tore (« à
vérifier »)\footnote{La page écrit « non coplanaire avec trois des quatre
pts », avec un mot incertain ; la précision sur le tore est de l'édition.}.

\pagerange{53}{55}
\subsection*{Les données d'une réalisation régulière (pages 53 à 55)}

Un repère $r_0 = (s_0, a_0, f_0)$ étant choisi, $\mathrm{G}$ est engendré par
$\sigma_0, \sigma_1, \sigma_2$, où $\sigma_0$ change le sommet, $\sigma_1$
l'arête, $\sigma_2$ la face. Une réalisation régulière équivaut à la donnée de
$\varphi_G : \mathrm{G} \to \operatorname{Aut}(E)$ --- c'est-à-dire de $\sigma_0,
\sigma_1, \sigma_2 \in \operatorname{Aut}(E)$ satisfaisant certaines conditions,
« qu'on explicitera » --- et du drapeau $\varphi_R(r_0) = (s_0, a_0, f_0) \in
\operatorname{Drap}_{012}(E)$ : comme $\mathrm{R}$ est un $\mathrm{G}$-torseur,
une application $\mathrm{G}$-équivariante $\varphi_R$ est déterminée par
$\varphi_R(r_0)$\footnote{La page écrit « $\varphi_R(r_0) \in r_0 = (s_0, a_0,
f_0)$ », peut-être pour « $= $ ». À la page 54, la deuxième ligne de la
définition des $\sigma_i$ porte $\sigma_1(a_0) = a_0$ en troisième position ;
c'est $\sigma_1(f_0) = f_0$.}. La condition pour que $\pi^{(i)}_E \circ
\varphi_R$ passe au quotient en $\operatorname{Fac}_i(X) \to
\operatorname{Fac}_i(E)$ ($i = 0, 1, 2$) est que le stabilisateur de la
$i$-facette de $r_0$ fixe son image :
\[
\sigma_0 a_0 = a_0,\ \sigma_0 f_0 = f_0 ; \qquad
\sigma_1 s_0 = s_0,\ \sigma_1 f_0 = f_0 ; \qquad
\sigma_2 s_0 = s_0,\ \sigma_2 a_0 = a_0 .
\]
C'est la condition a) de la page 6, pour $n = 2$.

\pagerange{56}{56}
\subsection*{Cartes homogènes (page 56)}

Sous « à vérifier », la page énonce, pour les « cartes homogènes de type $(p,
q)$ » --- les cartes standards $C_{p,q}$ du dossier 88, $p$ l'ordre des
sommets, $q$ celui des faces : \emph{$C_{p,q}$ est une géométrie d'incidence si
et seulement si $p, q \geq 2$, et un polyèdre combinatoire si et seulement si
$p, q \geq 3$.} Pour $p = 1$, une arête a ses deux extrémités confondues en
un sommet ; pour $p = 2$, deux arêtes ont les mêmes extrémités, et deux faces se
touchent le long de deux arêtes, ce que l'axiome de treillis du dossier 86
(page 2) interdit\footnote{La page écrit d'abord a) « est géométrie
d'incidence ssi $(p, q) = (1, 2)$ ou $(2, 1)$ » et b) « est polyèdre
combinatoire ssi $p = 2$ ou $q = 2$ », qui énumèrent les exceptions plutôt
qu'elles ne caractérisent ; puis, entre crochets, les énoncés qu'on donne. Elle
n'a pas biffé a) et b). L'explication par le dossier 86 est de l'édition.}.

\pagerange{58}{58}
\subsection*{Un feuillet isolé : groupes opérant sur un groupoïde (page 58)}

D'une autre encre, sans lien visible avec ce qui l'entoure : $\mathcal{C}$ un
groupoïde, $G \subset \operatorname{Aut}(\mathcal{C})$ transitif sur les objets,
$H = G_e$ le stabilisateur d'un objet $e$, $\operatorname{Ob}\mathcal{C} \simeq
G / H$ ; $\Gamma$ le normalisateur de $G$, $\Gamma^0$ son
centralisateur\footnote{La page écrit « centralisateur de $\mathcal{C}$ » ; on
lit « de $G$ ».}, avec $1 \to \Gamma^0 \to \Gamma \to \operatorname{Aut}(G)$ ;
un élément $\gamma$ de $\Gamma^0$ est connu sur les objets par $\gamma(e)$, d'où
une flèche $\Gamma^0 \to N(H)/H$, de noyau $\Gamma^1$, qui s'envoie dans
$\operatorname{Aut}_H(\pi_e)$, $\pi_e$ le groupe des automorphismes de $e$. La
même chose est esquissée pour un graphe $K$ sur lequel $G$ opère
transitivement sur les arêtes, $K_1 \simeq G / H_1 \to K_0 = G / H_0$. Le fait
sous-jacent est classique : les automorphismes d'un $G$-ensemble homogène $G/H$
forment le groupe $N_G(H)/H$. Les suites exactes du bas de la page, superposées
et en partie barrées, ne s'interprètent pas avec sûreté.

\pagerange{61}{62}
\section*{V. Le groupe modulaire et les solides de Platon (pages 61 et 62)}

La page 61 s'ouvre au milieu d'un calcul. Pour $n = 2$, soit $\Gamma^+_{p,q}$
le sous-groupe d'indice $2$ des rotations, engendré par $\rho_s$ et $\rho_f$, du
groupe $\Gamma_{p,q} = \mathfrak{G}_2 / \langle\langle \rho_s^p, \rho_f^q
\rangle\rangle$ ; avec $\sigma = \sigma_0\sigma_2$ (d'ordre $2$), on a $\rho_f\rho_s
= \sigma$ (convention ci-dessus). Quand l'un des ordres vaut $3$ et l'autre est
infini, ce groupe est le produit libre $\mathbf{Z}/2 * \mathbf{Z}/3$, et
\[
\Gamma^+_{3,\infty} \simeq \mathrm{SL}(2, \mathbf{Z}) / \{\pm 1\} =
\mathrm{PSL}(2, \mathbf{Z}),
\]
le groupe modulaire\footnote{La page note ce groupe $\Gamma^0_{\infty,3}$, dans
sa convention de la page 61 où $\rho_s = \sigma_0\sigma_1$ et $\rho_f =
\sigma_1\sigma_2$ ; par dualité, $\Gamma^+_{p,q} \simeq \Gamma^+_{q,p}$, et l'ordre
des indices n'importe pas pour les groupes. La première ligne, $1 \to R_n \to
\mathrm{SL}(2, \mathbf{Z})/\pm 1 \to \mathrm{SL}(2, \mathbf{Z}_2) \to 0$, avec un
indice $n$ douteux, dit que le noyau de la réduction modulo $2$ est engendré,
comme sous-groupe distingué, par le carré d'un élément parabolique.}. Pour le
voir, la page cherche les stabilisateurs des points fixes des homographies
$z \mapsto (az + b)/(cz + d)$ : un point fixe $z = x + iy$, $y \neq 0$, impose
$d = a - c\operatorname{Tr}(z)$, $b = -cN(z)$, et $\det = N(a - cz) = 1$. Pour $z =
i$, le stabilisateur est $\{\bigl(\begin{smallmatrix} a & -c \\ c & a
\end{smallmatrix}\bigr)\}$, et $\sigma = \bigl(\begin{smallmatrix} 0 & -1 \\ 1 &
0 \end{smallmatrix}\bigr)$, avec $\sigma^2 = -1$ ; pour $z = \zeta = e^{i\pi/3}$,
c'est $\{\bigl(\begin{smallmatrix} a & -c \\ c & a - c
\end{smallmatrix}\bigr)\}$, et l'élément d'ordre $3$ (modulo $\pm 1$) est
$\bigl(\begin{smallmatrix} 0 & 1 \\ -1 & 1 \end{smallmatrix}\bigr)$, de cube
$-1$. Leur produit est la translation $\bigl(\begin{smallmatrix} 1 & -1 \\ 0 & 1
\end{smallmatrix}\bigr)$, dont la puissance $n$-ième est $\bigl(\begin{smallmatrix}
1 & -n \\ 0 & 1 \end{smallmatrix}\bigr)$\footnote{La page écrit ce produit
$\rho'_s = \sigma\rho_f = \sigma_2\rho_f\sigma_2^{-1}$ ; dans ses notations,
$\sigma\rho_f = \sigma_0\sigma_2\sigma_1\sigma_2 = \sigma_2\rho_s\sigma_2^{-1}$,
conjugué de $\rho_s$, non de $\rho_f$. Les signes de la matrice de $\rho_f$ sont
repassés.}.

Réduisant modulo $n$, on obtient un épimorphisme
\[
\Gamma^+_{3,n} \longrightarrow \mathrm{SL}(2, \mathbf{Z}/n) / \{\pm 1\} \qquad (n
\geq 2),
\]
puisque l'élément parabolique devient d'ordre $n$ ; c'est un isomorphisme si et
seulement si $n \leq 5$, c'est-à-dire si et seulement si $\Gamma^+_{3,n}$ est
fini. La page dresse le tableau :
\[
\begin{array}{c|c|c|l}
n & \Gamma^+_{3,n} & \mathrm{SL}(2, \mathbf{Z}/n)/\pm 1 & \\ \hline
2 & \mathfrak{S}_3,\ 6 & 6 & \text{diédral d'ordre } 6 \\
3 & \mathfrak{A}_4,\ 12 & 12 & \text{tétraèdre} \\
4 & \mathfrak{S}_4,\ 24 & 24 & \text{cube (ou octaèdre)} \\
5 & \mathfrak{A}_5,\ 60 & 60 & \text{icosaèdre (ou dodécaèdre)} \\
6 & \mathbf{Z}^2 \rtimes \mathbf{Z}/6,\ \infty & 72 & \text{pavage du plan par des triangles}
\end{array}
\]
avec $\operatorname{card}\mathrm{SL}(2, \mathbf{F}_q) = q(q^2 - 1)$ et
$\operatorname{card}\mathrm{SL}(2, \mathbf{Z}/p^\nu) = \operatorname{card}
\mathrm{SL}(2, \mathbf{F}_p)\cdot p^{3(\nu - 1)}$. Tout cela est exact. C'est le
fait classique que le sous-groupe de congruence principal $\Gamma(n)$ est
engendré, comme sous-groupe distingué de $\mathrm{PSL}(2, \mathbf{Z})$, par la
puissance $n$-ième de la translation si et seulement si $n \leq 5$ ; pour $n
\geq 6$, le quotient $\Gamma^+_{3,n}$ est infini et la carte correspondante
n'est plus une carte de congruence\footnote{Fait classique, connu de Klein et
Fricke ; cité de mémoire. La dernière case du tableau, « pavage par triangles
orienté du plan », est très serrée ; « cas 3-diédral » pour $n = 2$ est une
lecture probable. Les noms de solides sont ceux de la page ; un solide et son
dual ont le même groupe.}.

\pagerange{64}{71}
\section*{VI. Les trois réflexions en coordonnées, et la quadrique invariante (pages 64 à 71)}

\pagerange{64}{66}
\subsection*{Les trois réflexions (pages 64 à 66)}

Pour $n = 2$, on refait le calcul des pages 9 à 18 directement. Le repère $r =
(s_0, a_0, f_0)$ donne $s_1 = \sigma_0 s_0$, $s_2 = \sigma_1 s_1$, $s_3 = \sigma_2
s_2$, $e_i = s_i - s_0$. Les points $t_2 = \sigma_0 s_2$ et $t_3 = \sigma_0 s_3$
sont dans les plans $(s_0, s_1, s_2)$ et $(s_0, s_1, s_3)$, $\sigma_2$ les échange
(puisque $\sigma_0\sigma_2 = \sigma_2\sigma_0$), et l'on écrit $t_2 = (\lambda,
\mu, 0)$, $t_3 = (\lambda, 0, \mu)$, d'où
\[
\sigma_0(x, y, z) = \bigl(1 - x + (\lambda - 1)(y + z),\ \mu y,\ \mu z\bigr), \qquad
\sigma_0^2 = 1 \iff \mu^2 = 1,\ (\lambda - 1)(\mu - 1) = 0 .
\]
De même $s'_3 = \sigma_1 s_3 = s_0 + \xi e_1 + \eta e_2 + (1 + \zeta)e_3$ donne
$\sigma_1(x, y, z) = (y + \xi z,\ x + \eta z,\ (1 + \zeta)z)$, et $\sigma_1^2 = 1$
équivaut à $(1 + \zeta)^2 = 1$ et $\xi + \eta + \xi\zeta = \xi + \eta + \eta\zeta = 0$.
Comme $\det\sigma_{1V} = -(1 + \zeta)$, et que le lieu fixe de $\sigma_1$ est un
plan si et seulement si $\zeta = 0$ et $\xi = -\eta$, on prend $\zeta = 0$,
$\xi = -\eta$\footnote{La page écrit la troisième coordonnée de $s'_3$ (son
$s''_3$) « $1 - \zeta$ » au lieu de $1 + \zeta$, et en tire « $\zeta^2 = 0$,
$\eta + \xi = \xi\zeta = \eta\zeta$ » ; avec $1 + \zeta$, l'involution donne
$\zeta(\zeta + 2) = 0$. La conclusion $\zeta = 0$, $\xi = -\eta$, qu'imposent le
déterminant $-1$ ou le plan fixe, est la même. Les points d'interrogation de la
page sous « $\mu = 1$ ? », « $\xi = -\eta$ ? », « $\zeta = 0$ ? » sont levés à la
page 66.}. Quant à $\sigma_0$, c'est une réflexion si et seulement si $\mu =
1$ ; pour $\mu = -1$ (et alors $\lambda = 1$), c'est la symétrie par rapport au
milieu de $s_0 s_1$\footnote{La ligne de la page 66 se lit « $\sigma_0$ est une
\emph{symétrie par rapport à un hyperplan} sauf si $\mu = 1$ », lecture
incertaine d'une ligne rapide ; ce qu'on donne est le calcul.}. Avec $\zeta =
0$, $\mu = 1$, $\lambda = 1 + \beta$, $\eta = 1 + \alpha$ :
\[
\begin{aligned}
\sigma_0(x, y, z) &= \bigl(1 - x + \beta(y + z),\ y,\ z\bigr), \\
\sigma_1(x, y, z) &= \bigl(y - (1 + \alpha)z,\ x + (1 + \alpha)z,\ z\bigr), \\
\sigma_2(x, y, z) &= (x, z, y),
\end{aligned}
\]
qui est la forme normale de la page 18 pour $n = 2$, avec $\beta = \alpha_0$,
$\alpha = \alpha_1$.

\pagerange{67}{69}
\subsection*{La fonction quadratique invariante (pages 67 à 69)}

$\rho_s = \sigma_1\sigma_2$ est d'ordre $p$ si et seulement si $\Psi_p(\alpha) =
0$, et $\rho_f = \sigma_0\sigma_1$ d'ordre $q$ si et seulement si $\Psi_q(\beta) =
0$ (page 67 ; c'est la proposition de la section I). On cherche les fonctions
quadratiques $f = ax^2 + by^2 + cz^2 + uyz + vzx + wxy + px + qy + rz + d$
invariantes par $\sigma_0, \sigma_1, \sigma_2$ ; on peut supposer $d = 0$. La
page trouve, à un facteur près,
\[
f(x, y, z) = x^2 + y^2 + z^2 - \gamma\, yz - \beta\, zx - \beta\, xy - (x + y + z),
\qquad \gamma = 2 + 2(\alpha + \beta) + \alpha\beta = (\alpha + 2)(\beta + 2) - 2 .
\]
Les fonctions quadratiques invariantes sont les combinaisons de $f$ et de $1$ ;
celles qui s'annulent en $s_0$ --- donc sur tous les sommets --- sont les
multiples de $f$\footnote{Vérifié sur machine : $f$ est invariante par les trois
générateurs pour $\alpha, \beta$ quelconques. L'unicité découle du calcul de
la page : l'invariance par $\sigma_0$ et $\sigma_1$ dans le plan $z = 0$ donne
$a = b = -p = -q$, $w = -\beta a$, celle par $\sigma_2$ donne $c = b$, $v = w$,
$r = q$, et $u$ est déterminé par $\sigma_1$.}. Sa partie quadratique $f_0$ a
pour matrice
\[
\begin{pmatrix} 2 & -\beta & -\beta \\ -\beta & 2 & -\gamma \\ -\beta & -\gamma & 2
\end{pmatrix}, \qquad
\delta' = \tfrac12 \det = -(\alpha + 2)(\alpha + \beta)(\beta + 2)^2 ,
\]
et $f$, homogénéisée, a pour discriminant $\delta(f) = (\alpha - 2)(\alpha +
2)(\beta + 2)^2$\footnote{La page écrit $\delta' = 4 - \beta^2(\gamma + 2) -
\gamma^2 = -(\alpha + 2)(\beta + 2)^2(\alpha - \beta)$, le dernier facteur lu avec
doute ; le calcul donne $(\alpha + \beta)$. $\delta(f)$ est exact, ainsi que
$(\gamma + 2)(\gamma - 4\beta - 6)$, l'autre forme qu'en donne la page. La
condition d'inversibilité est dite « $\alpha \neq \pm 2$, $\beta \neq -2$,
$\alpha + \beta \neq 0$ », avec des mots incertains.}. Si $\alpha + \beta$ est
inversible, le groupe a un unique point fixe, le \emph{centre}
\[
c = \Bigl(\frac{\alpha}{2(\alpha + \beta)},\ \frac{-1}{2(\alpha + \beta)},\
\frac{-1}{2(\alpha + \beta)}\Bigr), \qquad f(c) = \frac{2 - \alpha}{4(\alpha + \beta)} ,
\]
et l'antipodisme est la symétrie $x \mapsto 2c - x$, qui conserve $f$. $f(c)$
s'annule si et seulement si $\alpha = 2$ : la quadrique circonscrite est alors
un cône\footnote{Tout ceci vérifié sur machine. La page note le centre
$\sigma$, lettre déjà prise, et l'antipodisme d'un $a$ souligné.}. Puis vient
le tableau (page 69), dont toutes les valeurs ont été vérifiées :
\[
\begin{array}{l|c|c|c|c}
& \alpha\ (\text{sommet}) & \beta\ (\text{face}) & \delta' & \delta \\ \hline
\text{tétraèdre} & -1 & -1 & 2 & -3 \\
\text{cube} & -1 & 0 & 4 & -12 \\
\text{octaèdre} & 0 & -1 & 2 & -4 \\
\text{icosaèdre} & \alpha^2 + \alpha = 1 & -1 & 1 & -\alpha - 3,\ N = 5 \\
\text{dodécaèdre} & -1 & \beta^2 + \beta = 1 & \beta + 2,\ N = 1 & -9\beta - 15,\ N = 9 \\
\text{« pentagrammes »} & \alpha^2 + \alpha = 1 & \beta = \alpha' & \beta + 2,\ N = 1 & 4\alpha - 3,\ N = 5
\end{array}
\]
($N$ la norme de $\mathbf{Q}(\sqrt 5)$ à $\mathbf{Q}$). On en lit les
caractéristiques exceptionnelles : le centre est à l'infini si et seulement si
la caractéristique est $2$ (dans les six cas, $2(\alpha + \beta)$ est, au signe
près, une puissance de $2$ multipliée par une unité) ; $f_0$ est dégénérée en
caractéristique $2$ pour les trois premiers, jamais pour les trois derniers ;
$f$ est dégénérée en caractéristique $3$ (tétraèdre), $2$ et $3$ (cube), $2$
(octaèdre), $5$ (icosaèdre et pentagrammes), $3$ (dodécaèdre). Les
« pentagrammes » --- $\alpha$ et $\beta$ les deux racines distinctes de
$x^2 + x - 1$ --- sont les polyèdres réguliers étoilés de type
$\{5, 5/2\}$ et $\{5/2, 5\}$, le grand dodécaèdre et le petit dodécaèdre
étoilé de Kepler--Poinsot\footnote{Les noms sont les nôtres ; dans le symbole $\{q, p\}$, $q$ est la face et $p$ la figure de sommet. Les lignes de l'icosaèdre et du dodécaèdre, prises avec l'autre racine de $x^2 + x - 1$, donnent les deux autres, le grand icosaèdre $\{3, 5/2\}$ et le grand dodécaèdre étoilé $\{5/2, 3\}$. La page ajoute pour
chaque cas une ligne « $Q_E$ bisingulière ssi car.\ $\neq \ldots$ », où
« bisingulière » est lu avec doute et dont le sens n'est pas sûr ; on ne la
reprend pas. Le paragraphe qui clôt la page 69, marqué d'un « ? », est en
grande partie illisible.}.

\pagerange{70}{71}
\subsection*{Un cas hyperbolique (pages 70 et 71)}

La page 71, barrée en haut, dit ce qu'il faut pour une réalisation réelle qui
soit un vrai dessin : pour $p$, $q$ finis $\geq 3$, $\alpha$ et $\beta$ de la
forme $\zeta + \zeta^{-1}$, $\zeta = e^{2i\pi/m}$ ; $\alpha = -2$ ou $\beta = -2$
« donne des croisements » ; et pour $p = \infty$ (resp.\ $q = \infty$) il faut
$\alpha = 2$ (resp.\ $\beta = 2$). Elle traite $\beta = 2$, $\alpha \neq -2$ :
$\gamma = 2(2\alpha + 3)$, $\delta' = -16(\alpha + 2)^2$, et
\[
f_0(x, y, z) = (-x + y + z)^2 - 4(\alpha + 2)\,yz, \qquad f(c) = \frac{2 -
\alpha}{4(\alpha + 2)} .
\]
Pour $-2 < \alpha < 2$, $f(c) = \lambda > 0$, et, en coordonnées centrées $X =
x - c$, les sommets vérifient $f_0(X) = -\lambda$. Avec $4YZ = (Y + Z)^2 - (Y -
Z)^2$ et $U = (-X + Y + Z)/\sqrt\lambda$, $V = \sqrt{\alpha + 2}\,(Y -
Z)/\sqrt\lambda$, $W = \sqrt{\alpha + 2}\,(Y + Z)/\sqrt\lambda$, cela devient
\[
W^2 = U^2 + V^2 + 1 ,
\]
la dernière ligne de la page 70 : un hyperboloïde à deux nappes\footnote{Les
lettres $U, V, W$ sont surchargées sur la page, et plusieurs lignes
intermédiaires sont barrées ; on a refait le changement de variables pour qu'il
aboutisse à l'équation finale de la page.}. C'est le modèle de l'hyperboloïde
du plan hyperbolique, et le groupe, qui conserve la forme de signature $(2,
1)$, y agit par isométries : les faces y sont des « apeirogones » et l'on
obtient le pavage régulier $\{\infty, p\}$ du plan hyperbolique, avec $\alpha =
2\cos(2\pi/p)$ --- l'une des réalisations « géodésiques » que le dossier 88
attribue aux cartes standards\footnote{L'identification est de l'édition ; la
page s'arrête sur l'équation.}.

\pagerange{73}{111}
\section*{VII. L'icosaèdre en coordonnées universelles (pages 73 à 111)}

La chemise de la page 73 porte « (32) Calculs en coordonnées universelles
(icosaèdre) --- vieilles rédactions ». Dans toute la section, $\beta = -1$
(faces triangulaires) et $\alpha^2 + \alpha - 1 = 0$ (cinq faces en chaque
sommet), $\alpha' = -1 - \alpha$, de sorte que $\alpha\alpha' = -1$. Les calculs
ont lieu sur l'anneau $\mathbf{Z}[\alpha] = \mathbf{Z}[x]/(x^2 + x - 1)$, ou sur
toute base où $\alpha$ est défini. La page 89 note les identités utiles :
\[
(1 + \alpha)\alpha = 1, \quad (1 - \alpha)(1 - \alpha') = 1, \quad (1 + \alpha)(1 +
\alpha') = -1,
\]
\[
(\alpha + 2)(\alpha' + 2) = 1, \quad (2\alpha - 1)(2\alpha' - 1) = -1 ;
\]
en particulier $\alpha$, $1 \pm \alpha$, $\alpha + 2$, $2\alpha - 1$ sont des
unités, et $\alpha - 2$ en est une si et seulement si $5$ est inversible, car
$(\alpha - 2)(\alpha + 3) = -5$ (page 90 ; l'identité est de l'édition).

\pagerange{74}{75}
\subsection*{Première passe : les réflexions et les douze sommets (pages 74 et 75)}

Les trois générateurs sont
\[
\sigma_0(x, y, z) = \bigl(1 - (x + y + z),\ y,\ z\bigr), \quad
\sigma_1(x, y, z) = (y + \alpha' z,\ x - \alpha' z,\ z), \quad
\sigma_2(x, y, z) = (x, z, y),
\]
de miroirs $H_0 : 2x + y + z = 1$, $H_1 : x - y - \alpha' z = 0$, $H_2 : y = z$,
et de centres les directions $\Delta_0 = \langle e_1 \rangle$, $\Delta_1 =
\langle e_1 - e_2 \rangle$, $\Delta_2 = \langle e_2 - e_3 \rangle$. En
caractéristique $2$, $H_0$ devient $y + z = 1$ et $H_2$ s'écrit $y + z = 0$ :
ils sont parallèles\footnote{La note de la transcription doute que cela se lise
sur les équations transcrites ; cela s'y lit.}. Les intersections deux à deux
des miroirs sont les axes : $D_0 = H_1 \cap H_2$, la droite $y = z = \alpha' x$,
de direction $(1, \alpha', \alpha') = \alpha'(-\alpha, 1, 1)$, qui passe par
$s_0$ ; $D_1 = H_2 \cap H_0$, $y = z = \tfrac12 - x$ si $2$ est inversible ;
$D_2 = H_0 \cap H_1$, $y = \alpha + (1 - 3\alpha)x$, $z = (1 - \alpha)(1 - 3x)$
(page 111, vérifié). Le centre et l'antipodisme sont
\[
o = \tfrac12(\alpha',\ 1 - \alpha',\ 1 - \alpha'), \qquad
\iota(x, y, z) = \bigl(\alpha' - x,\ 1 - \alpha' - y,\ 1 - \alpha' - z\bigr) ,
\]
et les douze sommets
\[
\begin{array}{ll}
s_0 = (0, 0, 0) & \iota s_0 = (\alpha', 1 - \alpha', 1 - \alpha') \\
s_1 = (1, 0, 0) & \iota s_1 = (\alpha' - 1, 1 - \alpha', 1 - \alpha') \\
s_2 = (0, 1, 0) & \iota s_2 = (\alpha', -\alpha', 1 - \alpha') \\
s_3 = (0, 0, 1) & \iota s_3 = (\alpha', 1 - \alpha', -\alpha') \\
s'_3 = (\alpha', -\alpha', 1) & \iota s'_3 = (0, 1, -\alpha') \\
s''_3 = (0, -\alpha', 1) & \iota s''_3 = (\alpha', 1, -\alpha')
\end{array}
\]
sont l'orbite de $s_0$ sous le groupe engendré, d'ordre $120$\footnote{Vérifié
sur machine sur $\mathbf{Q}(\sqrt 5)$. La page 75 donne $\iota s''_3 = (\alpha',
1 - 2\alpha', -\alpha')$ ; le calcul donne $(\alpha', 1, -\alpha')$. Les restes de
calcul « $\tfrac12(-1 + \sqrt 5)$ » et « $\lambda = -2 + \sqrt 8$ » sous la
ligne de l'antipodisme ne se rattachent à rien.}.

La fonction quadratique de la page 68 devient
\[
q(x, y, z) = x^2 + y^2 + z^2 - \alpha\,yz + zx + xy - (x + y + z) ,
\]
de partie homogène $q_0$, de matrice $\bigl(\begin{smallmatrix} 2 & 1 & 1 \\ 1 & 2
& -\alpha \\ 1 & -\alpha & 2 \end{smallmatrix}\bigr)$ et de demi-discriminant
$\delta'(q_0) = 1$. Elle est normalisée par $q_0(e) = 1$ pour un vecteur arête
$e = s_1 - s_0$ ; pour un vecteur « coarête » $f = s_3 - s_2$, joignant deux
sommets à distance $2$, $q_0(f) = \alpha + 2 = \alpha'^2$. Dans les coordonnées
classiques, où les sommets sont les permutations circulaires de $(\pm x, \pm y,
0)$, cela donne $(y/x)^2 = \alpha'^2$, et, pour l'icosaèdre convexe réel ($\alpha
= \tfrac12(-1 + \sqrt 5)$), $y/x = -\alpha' = \tfrac12(1 + \sqrt 5)$ : le
rectangle d'or. En caractéristique $2$, la forme polaire de $q_0$ est
dégénérée, de noyau engendré par $u = (\alpha', \alpha, \alpha)$, mais $q_0(u) =
\alpha' \neq 0$, ce qui exprime que $q_0$ reste non dégénérée comme forme
quadratique\footnote{La page écrit le noyau $(\alpha', 1 - \alpha', 1 - \alpha')
= (\alpha', \alpha, \alpha) = (1 + \alpha, \alpha, \alpha)$, égalités valables en
caractéristique $2$. Le crochet de la page 75 récrit $q$ en caractéristique
$\neq 2$ sous une autre forme, que l'on ne reprend pas.}.

\pagerange{76}{78}
\subsection*{Vecteurs et directions remarquables ; la caractéristique 5 (pages 76 à 78)}

Le centre étant pris pour origine en caractéristique $\neq 2$, la forme
quadratique identifie $V$ à son dual, donc les directions de droites et de
plans. La page dresse la liste : $12$ vecteurs « doubles sommets » $S$ ; $30$
vecteurs arêtes $\Delta$ (quinze paires de vecteurs opposés) ; $30$ vecteurs
coarêtes $\Delta'$ ; les $30$ milieux d'arêtes $\mathcal{A}$ et de coarêtes
$\mathcal{A}'$ ; les $20$ barycentres (triplés) des faces $\Phi$ et des
« cofaces » $\Phi'$, triangles de sommets deux à deux à distance $2$ ; les
$20$ sommets $\Psi$ et $\Psi'$ de dodécaèdres circonscrits ; les $12$
barycentres $\Sigma$ des pentagones de sommets voisins. Elle écrit
\[
\Delta = \alpha\Delta', \quad \Delta = \mathcal{A}', \quad \Delta' = \mathcal{A},
\quad \Phi' = (2\alpha - 1)\Phi ,
\]
où $\mathcal{A}$, $\mathcal{A}'$ sont les sommes $s_i + s_j$ (centrées) ; ces
relations ont été vérifiées sur machine, avec $20$ cofaces. Les relations sur
$\Psi$ et $\Sigma$ --- « $\Psi = \mu\Phi$, $\mu = 1 + 2\alpha$, $\mu^2 = 5$ »,
« $\Sigma = \rho S$ » --- sont rapportées sans vérification : leurs légendes
sont trop incertaines pour qu'on sache quels points elles désignent\footnote{La
page est d'une écriture très rapide ; la note de la transcription dit les
listes « sûres dans leur structure, beaucoup moins dans leurs légendes ».
$\Phi$ et $\Phi'$ sont les sommets des dodécaèdres inscrits dans $\Pi$ et dans
le « cofacettoïde associé » $\Pi'$, lecture incertaine.}.

\textbf{Directions, et la caractéristique 5} (page 78). Il y a $6$ directions
de sommets, $15$ d'arêtes et $10$ de plans de faces : $31 = 5^2 + 5 + 1$, le
nombre de points du plan projectif sur $\mathbf{F}_5$. Ce n'est pas un hasard.
En caractéristique $5$, $\alpha^2 + \alpha - 1 = (\alpha - 2)^2$, donc $\alpha =
2$ et $f(o) = (2 - \alpha)/(4(\alpha - 1)) = 0$ : les sommets sont des vecteurs
isotropes de $q_0$, et les six directions de sommets sont les six points de la
conique $q_0 = 0$ sur $\mathbf{F}_5$. Les quinze directions d'arêtes
correspondent aux points extérieurs à la conique (intersections de deux
tangentes), les dix directions de faces aux points intérieurs, comme le dit la
page en termes de « sécantes » et de « diviseurs conjugués ». Pour les
vecteurs, $125 = 1 + 24 + 60 + 40$ : le vecteur nul, $24 = 2 \times 12$
vecteurs isotropes non nuls --- les sommets de « deux icosaèdres » ---, $60$
et $40$ vecteurs sur les points extérieurs et intérieurs\footnote{L'explication
par les points extérieurs et intérieurs est de l'édition ; c'est la géométrie
d'une conique sur $\mathbf{F}_q$, $q$ impair, qui a $q + 1$ points,
$q(q+1)/2$ points extérieurs et $q(q-1)/2$ points intérieurs. La page écrit
d'abord $1 + 12 + 60 + 40 = 113$, total surchargé. C'est l'isomorphisme
exceptionnel $\mathfrak{A}_5 \simeq \mathrm{PSL}(2, \mathbf{F}_5)$, que les
annonces de cours du dossier 75 mettent en tête du programme.}.

\pagerange{80}{80}
\subsection*{Un départ par les cocycles (page 80)}

Une représentation affine est une représentation linéaire $\Gamma \to
\mathrm{SO}(q)$ accompagnée d'un $1$-cocycle $f : \Gamma \to V$, $f(gg') = f(g)
+ g f(g')$. Avec $s_0$ pour origine, $f(\sigma_1) = f(\sigma_2) = 0$, et tout
revient à $e = f(\sigma_0)$ : $\sigma_0^2 = 1$ donne $\sigma_0 e = -e$, la
commutation de $\sigma_0$ et $\sigma_2$ donne $\sigma_2 e = e$, et, pour $u =
\rho_f$, $f(u^3) = (1 + u + u^2)e = 0$. La page s'arrête sur ces conditions
encadrées\footnote{La page est coupée à droite dans la numérisation. Elle
contient aussi une vérification de $\alpha' = -\alpha - 1$ et le calcul de $\bar
s_1 - \bar s_0 = (2\alpha - 1, 0, 0)$ pour des points $\bar s_0 = (-\alpha, 1,
1)$, $\bar s_1 = \sigma_0 \bar s_0$ de l'axe $D_0$ ; « $s_2 - s_0 = e_2 - e_3$ »
est tel quel.}.

\pagerange{82}{88}
\subsection*{La rédaction suivie : relations et nombre de sommets (pages 82 à 88)}

Les pages 82 à 99 reprennent tout, posément, en numérotant les formules. On
part de $\sigma_0$, $\sigma_2$ connus et de $\sigma_1$ inconnu, $\sigma_1 s_3 =
s'_3 = s_0 + \lambda e_1 + \mu e_2 + \nu e_3$. $\sigma_1^2 = 1$ s'écrit $\nu^2 =
1$, $\mu = -\nu\lambda$, $\lambda = -\nu\mu$ ; sur un corps, $\nu = 1$, $\mu =
-\lambda$ ou $\nu = -1$, $\mu = \lambda$ ; et $\det\sigma_1 = -\nu$, de sorte
qu'une réflexion demande $\nu = 1$. On pose $\lambda = -(1 + \alpha) = \alpha'$
(1), d'où
\[
s'_3 - s_3 = (1 + \alpha)(s_2 - s_1) \quad (2), \qquad
\sigma_1 = \begin{pmatrix} 0 & 1 & \alpha' & 0 \\ 1 & 0 & -\alpha' & 0 \\ 0 & 0 & 1
& 0 \\ 0 & 0 & 0 & 1 \end{pmatrix} \quad (3) .
\]
La page écrit la seconde forme de cette matrice avec un signe
douteux\footnote{Elle porte $-\alpha'$, $-\alpha'$ dans la troisième colonne,
le premier signe douteux ; avec $\alpha' = -(1 + \alpha)$, c'est $\alpha'$,
$-\alpha'$. Le signe de $\lambda(e_2 - e_1)$ dans le calcul de $s'_3 - s_3$ est
lu tel quel ; c'est $\lambda(e_1 - e_2)$.}.
Avec $\rho_f = \sigma_0\sigma_1$ et $\rho_s = \sigma_1\sigma_2$\footnote{La page
les appelle $u$ et $v$ ; elle écrit le cycle de $u$ sur les sommets du triangle
« $s_0 \mapsto s_1 \mapsto s_2 \mapsto s_1$ ? », pour $s_0$ en dernier.}, on
trouve $s''_3 = \rho_f s_3 = s_3 + (1 + \alpha)e_2$ (5), $s''_3 - s_3 = (1 +
\alpha)(s_2 - s_0)$ (6), et
\[
\rho_f : s_3 \mapsto s''_3 \mapsto s'_3 \mapsto s_3 \quad (7), \qquad \rho_f^3 =
\mathrm{id} \quad (8) .
\]
Les points $s_1, s_2, s_3, s'_3$ sont dans le plan $x + y + z = 1$, stable par
$\sigma_1$ et $\sigma_2$, donc par $\rho_s$, et « un calcul commode » montre
\[
\rho_s^5 = \mathrm{id} \iff \alpha^2 + \alpha - 1 = 0 \quad (9) ,
\]
ce qui est exact\footnote{C'est la proposition de la section I pour $p = 5$ ;
vérifié sur machine.}.
Le groupe $G = \langle \sigma_0, \sigma_1, \sigma_2 \mid \sigma_i^2,
(\sigma_0\sigma_2)^2, \rho_f^3, \rho_s^5 \rangle$ (10) --- le groupe de Coxeter
$H_3$, groupe des automorphismes de l'icosaèdre combinatoire épinglé --- opère
donc sur l'espace dès que $\alpha^2 + \alpha = 1$\footnote{Dans (10), la
relation de commutation se lit $\sigma_0\sigma_2 = \sigma_2\sigma_1$ ; c'est
$\sigma_2\sigma_0$. Le nom $H_3$ n'est pas sur la page.}.

Les transformés de $s_0$, de $\operatorname{Dr}(s_0, s_1)$ et de
$\operatorname{Pl}(s_0, s_1, s_2)$ forment-ils un icosaèdre géométrique ---
c'est-à-dire, sur tout corps, n'ont-ils entre eux que les incidences de
l'icosaèdre combinatoire ? (pages 85 à 87). « Si nous admettons que $G$ est
fini d'ordre $120$ » --- c'est vrai, et classique --- le stabilisateur de $s_0$
contient le groupe diédral $\langle \sigma_1, \rho_s \rangle$ d'ordre $10$, donc
le nombre de sommets divise $12$. Les six points $s_0, s_1, s_2, s_3, s'_3,
s''_3$ sont distincts : $s'_3$ et $s''_3$ ont pour dernière coordonnée $1$, et
ils coïncident entre eux ou avec $s_3$ si et seulement si $1 + \alpha = 0$, ce
qui, avec $\alpha^2 + \alpha = 1$, donnerait $0 = 1$ (« sur une base non vide
! »). Il y a donc $6$ ou $12$ sommets. S'il y en avait $6$, le stabilisateur de
$s_0$ serait d'ordre $20$ et contiendrait ce diédral ; or $G \simeq
\mathfrak{A}_5 \times \{1, \iota\}$, $\mathfrak{A}_5$ n'a pas de sous-groupe
d'ordre $20$, et le seul sous-groupe d'ordre $20$ qui contienne le diédral est
son produit par le centre $\{1, \iota\}$. Cela voudrait dire $\iota s_0 = s_0$
; le calcul des antipodes donne donc la liste des douze
sommets\footnote{L'argument est celui de la page 88, avec « il n'y en a qu'un
seul » ; la justification par $\mathfrak{A}_5$ est de l'édition. La page 86 ne
porte que la fin de phrase « de l'icosaèdre ». La page 81 n'est pas
transcrite.}.

\pagerange{89}{94}
\subsection*{Centre, antipodisme, incidences (pages 89 à 94)}

Le \emph{centre} est l'unique point fixe de $G$ (s'il existe) : sur les miroirs
(12)--(14), $o = \tfrac12(\alpha', 1 - \alpha', 1 - \alpha') = \tfrac12(-(1 +
\alpha), 2 + \alpha, 2 + \alpha)$ (15). En caractéristique $\neq 2$ c'est un
point ; en caractéristique $2$, c'est le point à l'infini dans la direction
$(-1, 1/\alpha, 1/\alpha)$.

\textbf{Proposition} (page 91). \emph{L'antipodisme $\iota$ --- l'élément central
non trivial de $G$, $(\sigma_0\sigma_1\sigma_2)^5$ --- est donné en toute
caractéristique par}
\[
\iota(x, y, z) = \bigl(\alpha' - x,\ (1 - \alpha') - y,\ (1 - \alpha') - z\bigr)
\quad (17) ;
\]
\emph{si $2$ est inversible, c'est la symétrie de centre $o$, et $o$ est son
unique point fixe ; en caractéristique $2$, c'est la translation de vecteur
$(\alpha', \alpha, \alpha)$, sans point fixe à distance finie.}\footnote{La
page n'identifie pas $\iota$ à un mot en les générateurs ; on a vérifié sur
machine que $(\sigma_0\sigma_1\sigma_2)^5$ est $x \mapsto 2o - x$. Elle lit le
vecteur de translation « $(\alpha', 1, 1)$ », les deux dernières composantes
surchargées ; c'est $(\alpha', 1 - \alpha', 1 - \alpha') = (\alpha', \alpha,
\alpha)$ en caractéristique $2$.} En caractéristique $2$, l'équation $\alpha^2 +
\alpha - 1 = 0$ n'a de solutions que dans $\mathbf{F}_4 \smallsetminus
\mathbf{F}_2$ : se donner $\alpha$ sur une base $Y$ au-dessus de $\mathbf{F}_2$,
c'est se donner une structure de $\mathbf{F}_4$-schéma, $Y \to \operatorname{Spec}
\mathbf{F}_4$, et $1/\alpha = \alpha + 1$ est l'autre élément de $\mathbf{F}_4
\smallsetminus \mathbf{F}_2$\footnote{La page note de la même lettre $K$ la
clôture algébrique et l'extension quadratique de $\mathbf{F}_2$ qu'elle
contient.}. Les antipodes des six premiers sommets (18) sont ceux du tableau
de la page 75 ci-dessus\footnote{Les deux dernières lignes de (18), page 92,
ont leur deuxième coordonnée noircie après le « $1$ » ; les valeurs exactes
sont $(0, 1, -\alpha')$ et $(\alpha', 1, -\alpha')$, et non celles qu'attend la
note de la transcription. La page 93 en donne d'abord une autre version,
barrée, calculée avec un autre centre.}.

\textbf{Incidences} (pages 93 et 94). La page veut prouver, « sur un corps de
base quelconque », que trois sommets distincts ne sont pas alignés, et que
quatre sommets distincts ne sont coplanaires que s'ils sont sur une même face
--- d'où il suivrait qu'un sommet ou une arête n'est incident géométriquement
à une face que s'il l'est combinatoirement. Une note marginale ajoute :
« c'est en fait \emph{faux} en car.\ $2$ ! ». La seconde affirmation est en
fait fausse en toute caractéristique : les cinq voisins d'un sommet sont
coplanaires (la page 85 le dit de quatre d'entre eux, $s_1, s_2, s_3, s'_3$),
et les extrémités de deux arêtes opposées aussi, sur un plan passant par le
centre --- le rectangle d'or ; il y a $75$ quadruplets coplanaires de sommets,
dont aucun ne contient une face. La page se replie sur ce qui suffit :
\begin{enumerate}
  \item[a)--c)] la droite qui joint $s_0$ à son antipode, à un voisin $s_1$, ou
    à un sommet $s''_3$ à distance $2$ ne contient pas d'autre sommet (« c'est
    vrai ») ; par homogénéité, aucun triplet de sommets n'est aligné ;
  \item aucun des neuf autres sommets n'est dans le plan $z = 0$ de la face
    $(s_0, s_1, s_2)$, c'est-à-dire n'a de coordonnée $z$ nulle --- « et on est
    heureux ! ».
\end{enumerate}
Ces deux faits valent en toute caractéristique, $2$ et $5$ comprises : les
coordonnées $z$ des neuf sommets sont $1$, $1 - \alpha' = 2 + \alpha$ ou
$-\alpha' = 1 + \alpha$, qui sont des unités de $\mathbf{Z}[\alpha]$ ; et pour
chaque triplet de sommets, les mineurs d'ordre $2$ qui expriment le non-alignement
engendrent l'idéal unité\footnote{Vérifié sur machine, pour tous les
triplets, toutes les paires (les douze points sont distincts en toute
caractéristique) et tous les couples (face, sommet) : les normes en jeu sont
$\pm 1$, ou de p.g.c.d.\ $1$. On a aussi contrôlé directement sur
$\mathbf{F}_4$, $\mathbf{F}_9$, $\mathbf{F}_5$, $\mathbf{F}_{49}$, $\mathbf{F}_{11}$
(deux racines), $\mathbf{F}_{169}$, $\mathbf{F}_{19}$, $\mathbf{F}_{29}$,
$\mathbf{F}_{31}$ : $12$ points distincts, aucun triplet aligné, $75$
quadruplets coplanaires, aucun contenant une face. Ce que vise la marge
« faux en car.\ $2$ » n'est pas clair ; en caractéristique $2$, $\iota$ étant
une translation, $s_0, s_1, \iota s_0, \iota s_1$ forment un parallélogramme,
mais ces quatre points sont déjà coplanaires en caractéristique $0$.}. La
réalisation est donc \emph{fidèle} au sens de la page 110 en toute
caractéristique : l'icosaèdre reste un icosaèdre.

\pagerange{95}{99}
\subsection*{La quadrique circonscrite (pages 95 à 99)}

Les fonctions quadratiques invariantes (page 95, puis page 97) : l'invariance
par $\sigma_2$ donne $b = c$, $v = w$, $\mu = \nu$ ; par $\sigma_0$, $v = a$ et
$\lambda = -a$ ; par $\sigma_1$ (formule (13)), $b = a$, $u = -\alpha a$, $\mu =
-a$. D'où
\[
\begin{aligned}
f &= a\,q, \\
q(x, y, z) &= x(x + y + z - 1) + (y^2 + z^2) - \alpha\,yz - (y + z) \\
&= x^2 + y^2 + z^2 - \alpha\,yz + zx + xy - (x + y + z) \qquad (19) .
\end{aligned}
\]
Les fonctions quadratiques invariantes sont les $a\,q + d$, et celles qui sont
nulles sur les sommets les $a\,q$\footnote{La page 95 encadre « $2a = 0$ ? »,
lecture douteuse, et garde un terme $\mu(y + z)$ qu'une ligne précédente
biffe ; la page 97 conclut proprement. La page 98, barrée, est une première
version du calcul de $f\sigma_1$, faite avec $\alpha$ là où il fallait
$\alpha'$, et aboutit à une autre forme, $x^2 + y^2 + z^2 - 2(1 + \alpha)yz - (x
+ y + z)$, que la page abandonne. La page 96, barrée aussi, commence le calcul
des quadriques passant par les sommets.}.

Les pages 99 et 101 font le calcul inverse, sans supposer l'invariance : une
fonction quadratique nulle en $s_0, s_1, s_2, s_3$ s'écrit $a(x^2 - x) + b(y^2 -
y) + c(z^2 - z) + uyz + vzx + wxy$ ; l'annulation en $s''_3$ donne $u = -\alpha
b$, puis, en $s'_3$, $b = c$ et $v = \alpha' w - (\alpha' - 1)a$, en $\iota s''_3$
donne $w = \alpha' v - (\alpha' - 1)a$, d'où $v = w = a$ (page 101), puis, en
$\iota s_2$, $b = a$. Donc \emph{les douze sommets sont sur une quadrique et une
seule, $q = 0$}, et la vérification en $\iota s_0$, $\iota s_1$ est inutile, $q$
étant invariante (« c'est clair »)\footnote{Vérifié sur machine : les
conditions d'annulation aux douze sommets sont de rang $9$ dans l'espace de
dimension $10$ des fonctions quadratiques. La page 101 encadre $w = a$ avec un
calcul intermédiaire en partie biffé.}.

\pagerange{100}{100}
\subsection*{Un feuillet étranger (page 100)}

Au crayon puis à l'encre, barré, un brouillon de topologie générale sans lien
avec l'icosaèdre : on construit une homotopie $h : X \times I \to Y$, égale à
$(f|F') \circ \mathrm{pr}_1$ sur $F' \times I$ et à $(b, t) \mapsto
c(t\varphi(b))$ sur $F \times I$, à l'aide d'une fonction continue égale à $1$
en un point $x$ et à $0$ sur le bord $B = F \cap F'$, avec $X = F \cup F'$ ;
une condition a) demande l'existence d'une telle fonction (« p.\ ex.\ si $F$
compact régulier », lecture douteuse). On n'en dit pas plus.

\pagerange{101}{106}
\subsection*{Discriminants ; deux feuillets intercalés (pages 101 à 106)}

La forme homogène $q_0$ a pour matrice $\bigl(\begin{smallmatrix} 2 & 1 & 1 \\ 1
& 2 & -\alpha \\ 1 & -\alpha & 2 \end{smallmatrix}\bigr)$, de déterminant $8 -
2\alpha - 2 - 2\alpha^2 - 2 = 4 - 2(\alpha^2 + \alpha) = 2$ : $\delta = 2$, donc
$\delta' = 1$ (20)--(21). « Forme quadratique très parfaite, \emph{merveille}
! » La forme $q$ homogénéisée par une variable $t$ a pour matrice
\[
\begin{pmatrix} 2 & 1 & 1 & -1 \\ 1 & 2 & -\alpha & -1 \\ 1 & -\alpha & 2 & -1 \\ -1
& -1 & -1 & 0 \end{pmatrix}, \qquad \det = -\alpha - 3, \quad N(-\alpha - 3) = 5 ,
\]
« sauf erreur de calcul » --- il n'y en a pas\footnote{La troisième ligne de
la matrice (22) est écrite $(-\alpha, -1, 2, -1)$, non symétrique ; on la
rétablit. La valeur $-\alpha - 3$ est celle du tableau de la page 69,
vérifiée. Le calcul de la page 104 par les mineurs $\Delta_1 = -\alpha'$,
$\Delta_2 = \alpha + 2$, $\Delta_3 = -2 - \alpha$, qui aboutit à $\alpha'$,
contient des erreurs de mineurs ; le résultat retenu à la page 102 est le
bon.}. La quadrique projective est donc lisse sauf en caractéristique $5$, où
--- ajout à l'encre noire --- « l'icosaèdre est inscrit dans une \ldots\ de
rayon nul, des quadriques \emph{sing.}\ ! » : le cône de sommet $o$ que l'on a
vu à la page 78. En toute autre caractéristique, l'icosaèdre est inscrit dans
une « quadrique \emph{lisse} » (fin de la phrase, en tête de la page
106)\footnote{Une première version, en bleu, biffée, disait « Vérifier ! Donc
l'icosaèdre est inscrit dans une quadrique lisse en toute dimension ! » ; elle
se poursuit page 106. Le mot après « 5- » n'est pas lu.}.

Deux feuillets d'un autre papier sont intercalés. La page 103 considère le
groupe $\mu_n$ opérant sur un module quadratique de rang $2$, $V_T \simeq D \oplus
D^{-1}$, par $\zeta.(x, y) = (\zeta x, \zeta^{-1}y)$ ; elle pose $c = \zeta +
\zeta^{-1}$, $\zeta^2 - c\zeta + 1 = 0$, et note qu'une section $c$ de
$\mathcal{O}_S$ racine du polynôme cyclotomique « réel » définit un revêtement
quadratique $\operatorname{Spec}\mathcal{O}_S[X]/(X^2 - cX + 1)$ sur lequel vit
$\zeta$, avec $\sigma\zeta = \zeta^{-1}$ pour l'involution du revêtement. C'est,
dans la langue des schémas, le passage du paramètre $\alpha = \zeta +
\zeta^{-1}$ des rotations à la racine de l'unité $\zeta$ elle-même\footnote{La
page écrit « $c$ satisfaisant $\Phi_n(c) = 0$ » ; c'est le polynôme qu'on a noté
$\Psi_n$ (pour $\Phi_n$ cyclotomique, $\Phi_n(\zeta) = 0$). Le rapprochement
avec le paramètre $\alpha$ est de l'édition. Une liste numérotée 1) à 3) des
données équivalentes (module quadratique de rang $2$, section de
$\Sigma^{(q)}/\mu_n^{(q)}$, revêtement étale quadratique $\Omega$ et
isomorphisme) est en partie illisible.}. La page 105 note que pour un module
localement libre $F$ de rang $2$, $\mathbb{V} = \operatorname{Sym}^2(F) \otimes
(\det F)^{-1}$ porte une forme quadratique canonique, le discriminant $Q :
\mathbb{V} \to \mathcal{O}_S$, qui en base s'écrit $b^2 - ac$ : c'est la
construction de l'isomorphisme $\mathrm{PGL}_2 \simeq \mathrm{SO}_3$ que la
lecture du dossier 82 développe\footnote{Feuille de listing, plusieurs lignes
barrées ; on ne donne que ce qui se lit. La normalisation exacte (« $-xx' +
b(xy' + x'y) + c\,yy'$ ») est incertaine.}. La page 104 calcule les mineurs du
déterminant (22).

\pagerange{107}{111}
\subsection*{Suffisance ; réalisations régulières et fidèles (pages 107 à 111)}

« Est-il clair que la condition $\alpha^2 + \alpha - 1 = 0$ est \emph{suffisante}
pour avoir une représentation de l'icosaèdre ? » (page 107). On a une
représentation affine du groupe de Coxeter $\Gamma' = \langle \sigma_0,
\sigma_1, \sigma_2 \mid \sigma_i^2, (\sigma_0\sigma_1)^3, (\sigma_1\sigma_2)^5,
(\sigma_0\sigma_2)^2\rangle$, qui a pour quotient le groupe $\Gamma$ de
l'icosaèdre combinatoire ; si l'on sait que $\Gamma' \simeq \Gamma$, c'est
fini, sinon il faut vérifier, « de pure routine », que les douze points sont
stables par les générateurs et que ceux-ci y opèrent comme sur les douze
sommets de l'icosaèdre combinatoire. On sait aujourd'hui que $\Gamma' \simeq
\Gamma$ (les deux sont d'ordre $120$) ; et la vérification de routine a été
faite ici sur machine\footnote{La page 108 résume : pour $\sigma_1$, des
« modules » $\lambda, \mu, \nu$ avec $\nu^2 = 1$, $\mu = -\nu\lambda$, puis $\nu = 1$
« exprimant $u^3 = \mathrm{id}$ » et $\lambda^2 + \lambda - 1 = 0$. Ce n'est pas
$\rho_f^3 = \mathrm{id}$, automatique (page 85), mais $\rho_s^5 = \mathrm{id}$
qui donne l'équation ; $\lambda = \alpha'$ vérifie la même équation que $\alpha$.
Le dernier paragraphe de la page 108, barré, voulait traiter le cas $\nu = -1$.}.

\textbf{Définitions} (pages 109 et 110). Soit $\Pi = (S, A, F; r)$ l'icosaèdre
combinatoire épinglé par le repère $r = (s_0, \operatorname{Dr}(s_0, s_1),
\operatorname{Pl}(s_0, s_1, s_2))$, et $\Gamma$ son groupe d'automorphismes. Une
\emph{représentation} (ou réalisation) affine de $\Pi$ dans un fibré affine $E$
de dimension relative $3$ sur un schéma $Y$ est la donnée d'applications
$\rho_S$, $\rho_A$, $\rho_F$ de $S$, $A$, $F$ dans les sections de $E$, ses
sous-fibrés en droites affines, ses sous-fibrés en plans affines, qui
respectent l'incidence\footnote{La page note $S$ la base, comme les
sommets ; on l'appelle $Y$.}. Elle est
\begin{itemize}
  \item \emph{acceptable} si les sommets d'une face ont des images deux à deux
    distinctes --- alors $\rho_A$, $\rho_F$ sont connues par $\rho_S$, qui
    peut être donnée arbitrairement ;
  \item \emph{fidèle} si deux facettes dont les réalisations sont incidentes
    sur une fibre sont incidentes ; cela entraîne qu'elle est acceptable et que
    $\rho_S$, $\rho_A$, $\rho_F$ sont injectives ;
  \item \emph{régulière} si elle est acceptable et si tout $g \in \Gamma$ se
    prolonge en un automorphisme affine $g_E$ de $E$ ; c'est une propriété de
    $\rho_S$, qui entraîne que $\rho_S(S)$ engendre affinement $E$.
\end{itemize}

\textbf{Proposition} (page 109). \emph{Une représentation régulière est fidèle.}
La démonstration commence comme celle de la page 51 --- $s_0, s_1, s_2, s_3$ sont
affinement indépendants, car sinon le plan $P = \operatorname{Pl}(s_0, s_1, s_2)$,
déjà stable par $\sigma_0$ et $\sigma_1$, le serait aussi par $\sigma_2$ --- et
s'interrompt. Pour l'icosaèdre, la conclusion est vraie, et c'est le résultat
des pages 93 et 94 complété plus haut : la seule réalisation régulière à repère
près (celle des pages 82 à 99) est fidèle en toute caractéristique.

\textbf{Réalisations régulières et représentations} (page 111). Se donner une
réalisation régulière de $\Pi$ revient à se donner une représentation affine de
rang $3$ de $\Gamma$ dans $E$ et un point $s_0 \in E^{\sigma_1} \cap
E^{\sigma_2}$, c'est-à-dire sur l'axe $D_0$, tel que $s_0$, $s_1 = \sigma_0 s_0$,
$s_2 = \sigma_1 s_1$, $s_3 = \sigma_2 s_2$ soient affinement indépendants. Deux
tels points $s_0, s'_0$ se correspondent par un unique automorphisme de $E$
commutant à $\Gamma$ --- l'homothétie de centre $o$ qui envoie l'un sur
l'autre, quand $2$ est inversible\footnote{L'homothétie est de l'édition : la
représentation linéaire est absolument irréductible, son commutant est formé
des scalaires, et un automorphisme commutant à $\Gamma$ fixe $o$. La marge de
la page 111, en partie illisible, esquisse un programme : représentations
affines fidèles, conditions d'existence, « classification des représentations
affines \ldots\ des formes de l'icosaèdre ».}.

\pagerange{113}{119}
\section*{VIII. Position du problème des 2-polyèdres réguliers (pages 113 à 119)}

La chemise rose de la page 113 porte « Position du pb des 2-polyèdres réguliers
/ vieille rédaction » et « (10 p) ».

\pagerange{114}{115}
\subsection*{Drapeaux universels (pages 114 et 115)}

Soit $\mathfrak{G}_2 = \langle \sigma_0, \sigma_1, \sigma_2 \mid \sigma_i^2,
(\sigma_0\sigma_2)^2 \rangle$ (1)--(2) le groupe cartographique\footnote{La page
le note $\Gamma$.}. Pour toute partie non vide $H \subset \{0, 1, 2\}$, soit
$\mathfrak{G}_H$ le sous-groupe engendré par les $\sigma_i$, $i \notin H$, et
$\mathfrak{F}_H = \mathfrak{G}_2 / \mathfrak{G}_H$ l'ensemble des \emph{drapeaux
universels de type $H$} (3). Pour $H \subset K$, $\mathfrak{G}_K \subset
\mathfrak{G}_H$ donne $\varphi_{HK} : \mathfrak{F}_K \to \mathfrak{F}_H$ (4), et
$\mathfrak{F}_{**} = \coprod_H \mathfrak{F}_H$ est ordonné par $x \preccurlyeq y$
si et seulement si $H \subset K$ et $x = \varphi_{HK}(y)$ (5)--(6). Les éléments
minimaux, $\mathfrak{F}_0$, $\mathfrak{F}_1$, $\mathfrak{F}_2$, sont les
sommets, arêtes et faces universels ; sur leur réunion $\mathfrak{F}_*$, on a la
dimension (8) et la relation d'\emph{incidence} : $x \leq y$ si $\dim x \leq
\dim y$ et s'il existe $z \in \mathfrak{F}_{**}$ avec $x, y \preccurlyeq z$
(9)\footnote{La page écrit $z \in \mathfrak{F}_*$ ; on lit
$\mathfrak{F}_{**}$, comme le demande le sens. Pour que ce soit une relation
d'ordre, il faut que l'incidence soit transitive, ce qui vaut pour la carte
universelle, et que la page tient pour acquis.}. C'est la carte universelle,
dont toute $2$-carte est un quotient.

\pagerange{115}{117}
\subsection*{Un polyèdre régulier dans un schéma (pages 115 à 117)}

Soit $X$ un schéma affine (ou quasi-projectif) sur une base $Y$, et
$\mathfrak{F}_H(X/Y)$ ses drapeaux de type $H$, au sens habituel. Un
\emph{polyèdre régulier scindé} $P$ dans $X$ est la donnée de trois ensembles
finis
\[
S \subset \mathfrak{F}_0(X/Y) = X(Y), \qquad A \subset \mathfrak{F}_1(X/Y), \qquad F
\subset \mathfrak{F}_2(X/Y) \quad (10)
\]
satisfaisant [11] : (a) chaque arête est incidente à exactement deux
sommets, (b) chaque couple sommet--face incidents l'est à exactement deux
arêtes, (c) chaque arête à exactement deux faces --- et chaque condition reste
vraie après tout changement de base $Y' \to Y$, $Y' \neq \varnothing$. Alors
$\Psi_* = S \amalg A \amalg F$ (12) est un ensemble ordonné du type voulu, on en
forme les drapeaux de tout type $H$ (12 bis), et $\operatorname{Rep}(P) =
\Psi_{\{0,1,2\}}$ (13), sur lequel $\mathfrak{G}_2$ opère par (14), $\sigma_i$
changeant la $i$-ième composante. Les restrictions $\operatorname{Rep}(P) \to
\Psi_H(P)$ (15) se factorisent par $\operatorname{Rep}(P)/\mathfrak{G}_H$ (16). On
demande de plus :
\begin{enumerate}
  \item[{[17]}] les applications (15) sont surjectives ;
  \item[{[18]}] $S$ engendre affinement $X$ fibre par fibre ;
  \item[{[19]}] une condition de connexité sur $\Psi_*(P)$, qui assure que
    $\mathfrak{G}_2$ opère transitivement sur $\operatorname{Rep}(P)$ et qu'un
    automorphisme de $P$ fixant un repère est l'identité ;
  \item[{[20]}] $\operatorname{Aut}(X, P)$ opère transitivement sur
    $\operatorname{Rep}(P)$\footnote{En marge de [19], « redondant ? », avec une
    condition équivalente proposée, $\operatorname{Rep}(P) / \mathfrak{G}_H
    \xrightarrow{\sim} \Psi_H(P)$ pour tout $H$. Le texte de [17] est en partie
    illisible au-delà de la première ligne.}.
\end{enumerate}
Alors le stabilisateur $N$ d'un repère est distingué --- les automorphismes de
$P$ commutent à $\mathfrak{G}_2$ et sont transitifs, de sorte que tous les
repères ont le même stabilisateur --- et $\operatorname{Rep}(P)$ est un torseur sous $\Gamma_P =
\mathfrak{G}_2 / N$ (21) ; $\operatorname{Aut}(X, P) \simeq
\operatorname{Aut}_{\mathrm{comb}}(P)$ (22), et le choix d'un repère $r$ définit
une bijection $\varphi : \Gamma_P \to \operatorname{Aut}(X, P)$, $\varphi(\gamma)$
étant l'unique automorphisme tel que $\varphi(\gamma)(r) = \gamma.r$ (23)--(24).
Ce $\varphi$ est un \emph{anti}-isomorphisme : $\varphi(\gamma\gamma')(r) =
\gamma\gamma'.r$, tandis que $\varphi(\gamma)\varphi(\gamma')(r) =
\varphi(\gamma)(\gamma'.r) = \gamma'.\varphi(\gamma)(r) = \gamma'\gamma.r$\footnote{La
page écrit « isomorphisme ». La même inversion de l'ordre est relevée par la
lecture du dossier 86, page 14.}.

\pagerange{117}{119}
\subsection*{Questions (pages 117 à 119)}

$r$ étant fixé, $\gamma \mapsto \gamma.r$ est surjective par [19], et les
surjections induites $N \backslash \mathfrak{F}_H = \mathfrak{G}_2 / N\mathfrak{G}_H
\to \Psi_H(P)$ (26) sont-elles bijectives ? Que dire des sous-groupes
distingués $N$ et des espaces homogènes $\mathfrak{G}_2 / N \simeq \Psi_{**}(P)$ ?
« On voudrait que ce soit l'ensemble des $5$ cartes correspondant aux polyèdres
\emph{réguliers} connus --- en particulier, on voudrait que ce soit
\emph{orientable} ! Et il faudrait, pour une carte combinatoire donnée,
déterminer les polyèdres réguliers épinglés par ce $\mathfrak{F}$ --- ce qui est
un pb modulaire\ldots »\footnote{« orientable » est une lecture incertaine.}.

Si (26) est bijective --- et sinon, il faudrait renforcer la définition --- on
récupère $P$ à partir de $\varphi : \mathfrak{G}_2 \to \operatorname{Aut}(X)$ (27) et
de $r = (s_0, a_0, f_0)$ (28) satisfaisant les conditions de stabilisateur (29),
qui sont celles de la page 55, en prenant pour $S, A, F$ les orbites. Pour
que de telles données définissent un polyèdre régulier (sur un corps), il faut
que $N = \ker\varphi$ (30) soit d'indice fini ; [17] est alors automatique,
mais [11] demande, avec $s'_0 = \sigma_0 s_0$, $a'_0 = \sigma_1 a_0$, $f'_0 =
\sigma_2 f_0$ (31), distincts de $s_0, a_0, f_0$ :
\begin{enumerate}
  \item[{[32]}] a) si $\gamma s_0 \prec a_0$, alors $\gamma s_0 \in \{s_0, s'_0\}$
    ; b) si $s_0 \prec \gamma a_0 \prec f_0$, alors $\gamma a_0 \in \{a_0,
    a'_0\}$ ; c) si $a_0 \prec \gamma f_0$, alors $\gamma f_0 \in \{f_0, f'_0\}$ ;
  \item[{[33]}] les $\gamma s_0$ engendrent affinement $X$ ;
  \item[{[33']}] si $\gamma s_0 \prec \gamma' a_0 \prec \gamma'' f_0$, il existe
    $\gamma'''$ avec $\gamma''' s_0 = \gamma s_0$, $\gamma''' a_0 = \gamma' a_0$,
    $\gamma''' f_0 = \gamma'' f_0$ --- ce qui exprime [19].
\end{enumerate}
Ce sont, pour l'icosaèdre, exactement les vérifications des pages 85 à 94 :
[32] et [33'] sont des énoncés d'incidence géométrique, et la fidélité établie
plus haut les contient\footnote{Le rapprochement est de l'édition. La page
renvoie à « [12] » pour [33], numéro surchargé ; la condition est [18].}.

\pagerange{120}{120}
\section*{IX. Les « cocubes » (page 120)}

D'une encre noire et d'une écriture posée. Soit $B \to I$ un revêtement double
d'ensembles finis, c'est-à-dire un ensemble $B$ muni d'une involution sans point
fixe $b \mapsto \bar b$, $I = B/\sigma$, $n = \operatorname{card} I$. Dans $k^B$,
soit $E$ l'espace affine des $\varphi = (\varphi_b)$ tels que $\varphi_b +
\varphi_{\bar b} = 1$ pour tout $b$ --- l'image réciproque de $(1, \dots, 1)$
par la trace $k^B \to k^I$ ---, de dimension $n$, et $H_b = \{\varphi_b = 1\}$.
Pour une section partielle $B' \in \Gamma_p(B/I)$, $F_{B'} = \bigcap_{b \in B'}
H_b$ ; l'ensemble $\mathfrak{F} = \{F_{B'}\} \cup \{\varnothing\}$ est celui des
faces du cube $[0, 1]^I$ : les sections totales donnent les sommets, la section
vide l'espace entier, et $H_b \cap H_{\bar b} = \varnothing$. C'est le cube
d'un ensemble à involution du dossier 76 (pages 49 et 50), où les faces sont
indexées, en renversant l'ordre, par les sections partielles\footnote{Le
dossier 76 prend $\varphi(\bar b) = -\varphi(b)$, d'où le cube $[-1, 1]^I$ ; ici
$\varphi_b + \varphi_{\bar b} = 1$ donne $[0, 1]^I$.}.

Dualement, soit $N \subset k^B$ le sous-module, facteur direct de rang $n - 1$,
engendré par les $(b + \bar b) - (b' + \bar b')$, et $V = k^B / N$, de rang $n +
1$, avec $0 \to W \to V \to k \to 0$. Les images des $b$ sont $2n$ points d'un
hyperplan affine de $V$ ; c'est le \emph{cocube}, que la page caractérise ainsi.
Un ensemble $\mathfrak{F}$ de sous-espaces affines d'un espace $E$ de dimension
$n$ est un cocube si
\begin{enumerate}
  \item[a)] pour tout point $s \in \mathfrak{F}_0$, il existe un unique $s' \neq
    s$ dans $\mathfrak{F}_0$ tel que la droite $ss'$ ne soit pas dans
    $\mathfrak{F}$ ; $s \mapsto s'$ est une involution sans point fixe de $B =
    \mathfrak{F}_0$, et $I = B/\sigma$ ;
  \item[b)] pour toute section partielle $B'$, $B'$ est affinement libre et son
    enveloppe affine est dans $\mathfrak{F}$, et l'on obtient ainsi tous les
    éléments de $\mathfrak{F}$ autres que $E$ ;
  \item[c)] pour toute section totale $B'$ et tout $b \notin B'$, $B' \cup
    \{b\}$ est affinement libre et engendre $E$ ;
  \item[d)] $b + \bar b = b' + \bar b'$ pour tous $b, b'$ : les milieux des
    diagonales coïncident.
\end{enumerate}
C'est le \emph{polytope croisé} --- l'octaèdre en dimension $3$ ---, dont les
faces sont les simplexes sur les ensembles de sommets sans paire opposée ; le
cube et lui sont duaux, et le dossier 87 les traite par la polarité
convexe\footnote{Les noms « polytope croisé », « hyperoctaèdre » ne sont pas sur
la page. Elle écrit « de dim.\ $n$ » après « sous-var.\ linéaires affines »,
pour l'espace $E$.}.

\pagerange{121}{123}
\section*{X. Programme provisoire d'une théorie arithmétique des polyèdres réguliers (pages 121 à 123)}

Le titre, souligné, est le sien, et la page commence un texte neuf.

Soit $\tau$ un type de polyèdre régulier dans l'espace affine réel de dimension
$n$. Il a un représentant « canonique » $X_\tau$, défini à isomorphisme
canonique près par le choix d'un repère $R_\tau$ ; l'espace $E_\tau$ qui le
porte a pour origine le barycentre des sommets et devient un espace vectoriel.
Soit $\mathcal{F}_\tau$ l'ensemble ordonné des facettes de $X_\tau$, y compris
celles de dimension $-1$ et $n$, et
\[
G_\tau \simeq \operatorname{Aut}_{\mathrm{aff}}(X_\tau) \simeq \operatorname{Aut}_{\mathrm{ord}}(\mathcal{F}_\tau) .
\]
La catégorie des polyèdres réguliers de type $\tau$ équivaut à celle des
torseurs sous $G_\tau$ --- par $X \mapsto \operatorname{Isom}(X_\tau, X)$\footnote{Le
foncteur est de l'édition ; la page énonce l'équivalence. L'indice du second
$\operatorname{Aut}$ se lit « ord » ou « ad ».}.

Soit $k$ un anneau commutatif (« plus gén., on pourrait se placer sur un schéma,
ou sur un topos loc.\ annelé »). Une \emph{$k$-réalisation} de $\tau$ est la
donnée
\begin{enumerate}
  \item[a)] d'un $k$-espace projectif $P$ de dimension $n$, « splitté » ---
    associé à un module projectif de type fini, défini à tensorisation près par
    un module inversible ---, muni d'une section $s_0$ et d'un hyperplan « à
    l'infini » $H_\infty$ ;
  \item[b)] d'un ensemble $\mathcal{F}$ de sous-espaces projectifs de $P$ tel
    que :
\end{enumerate}
\begin{enumerate}
      \item[1)] $(\mathcal{F}, \subset)$ soit isomorphe à $\mathcal{F}_\tau$, la
        dimension combinatoire de chaque $V \in \mathcal{F}$ (le plus petit
        élément étant de dimension $-1$) étant égale à sa dimension projective,
        et $P \in \mathcal{F}$ ;
      \item[2)] pour tout $V \in \mathcal{F}$, $s_0 \notin V$ et $V \not\subset
        H_\infty$, fibre par fibre ;
      \item[3)] les éléments de $\mathcal{F}_0$ engendrent $P$, et ceux de
        $\mathcal{F}_{n-1}$ engendrent le dual $\check P$ ; plus précisément,
        pour $V \subsetneq W$ dans $\mathcal{F}$, $W$ est engendré par les $V' \in
        \mathcal{F}$ avec $V \subset V' \subset W$, $\dim V' = \dim V + 1$, et
        dualement ;
      \item[4)] tout automorphisme de l'ensemble ordonné $\mathcal{F}$ est
        induit par un automorphisme de $P$ qui conserve $\mathcal{F}$.
\end{enumerate}
« Raisonnable si $\operatorname{Spec} k$ connexe ».\footnote{La condition 2)
est entre crochets, son numéro noirci, et une flèche la contourne : elle est
peut-être retirée. La précision de 3) est en marge et remplace une parenthèse
barrée. Dans 1), la page écrit « la dimension combinatoire de $\mathcal{F}$ »,
pour $V$.} C'est le renversement annoncé : on ne part plus d'une carte
combinatoire pour chercher ses réalisations, mais du polyèdre réel, et l'on
demande sur quels anneaux il « s'écrit » --- la question que les pages 74 à 111
résolvent pour l'icosaèdre, $\mathbf{Z}[\alpha]/(\alpha^2 + \alpha - 1)$ étant
l'anneau où il vit, avec $s_0$ le centre et $H_\infty$ l'hyperplan à l'infini
qui, en caractéristique $2$, contient ce centre. Aucun énoncé ne suit la
définition.

La page 123 est une feuille de notations, sans texte suivi, avec un cube et un
polygone où un drapeau est repassé en gras, et la chaîne $F_0, \dots, F_{n-1}$
d'un drapeau. On y lit des suites $(s_i)$, $(s'_i)$, $(\delta_i)$, $(f_i)$ avec
$s_{n-1} = 2$, $s_n = 1$, $s'_0 = 2$, $s'_{-1} = 1$, $f_{-1} = f_n = 1$, la
relation $s'_i f_{i+1} = f_i s_{i+1}$ et $N = \operatorname{card} G = \prod s_i =
\prod s'_i$. La lecture la plus naturelle : $f_i$ est le nombre de
$i$-facettes, $s_i$ le nombre de $i$-facettes contenant une $(i-1)$-facette
donnée, $s'_i$ le nombre de $i$-facettes contenues dans une $(i+1)$-facette
donnée ; la relation compte de deux façons les couples incidents de dimensions
$i$, $i+1$, et le nombre de drapeaux, qui est l'ordre du groupe, est
$s_0 s_1 \cdots s_n$ (pour le cube, $8 \cdot 3 \cdot 2 \cdot 1 = 48$)\footnote{Cette
interprétation est de l'édition ; elle rend compte des valeurs aux bords. Les
bornes du produit se lisent $1$ et $n - 1$, qui ne donnent pas l'ordre du
groupe. Les $\delta_i$ ($0 \leq i \leq n - 2$) ne sont pas définis ; ce sont
peut-être les $p_i$ de la page 3.}.

\end{document}
